\begingroup
\section{Node JP: exact support Hessian and the closed disk endpoint}
\label{module:JP}
\noindent\textit{Audited source: \nolinkurl{JP_joint_principal.tex}.}
\subsection{Fixed-normal support chart}

Fix cyclic unit normals $\nu_i$ with exterior angles
$\alpha_i\in(0,\pi)$ and $\sum_i\alpha_i=2\pi$.  For a support vector
$h=(h_i)$ put
\begin{equation}\label{JP:eq:P}
 P_\alpha(h)=\{x\in\mathbb R^2:x\cdot\nu_i<h_i\ \hbox{for all }i\},
 \qquad
 \Phi_\alpha(h)=\frac{|P_\alpha(h)|}{\pi}\lambda_1(P_\alpha(h)).
\end{equation}
We work in an active star on which all sides persist.  Translation vectors
$(b\cdot\nu_i)_i$ are quotiented exactly.  Although $\Phi_\alpha$ is also
dilation invariant, at a nonstationary point its Hessian does not descend
to an abstract dilation quotient.  We therefore fix an affine scale section
below.

Let $p_i(h)$ be the intersection of the two support lines with normals
$\nu_i,\nu_{i+1}$.  Solving the two linear equations shows that
\begin{equation}\label{JP:eq:vertex-linear}
 p_i(h+tz)=p_i(h)+t p_i(z).
\end{equation}
Consequently the physical-affine material parametrization of every side is
affine in $t$.  Its normal velocity on side $S_i(t)$ is exactly
\begin{equation}\label{JP:eq:constant-normal}
 V(t)\cdot\nu_i=z_i,\qquad \partial_\tau(V\cdot\nu_i)=0.
\end{equation}
This is the regular finite-energy replacement for the discontinuous
piecewise-constant function obtained by writing $z_i$ in the radial angular
coordinate.  That raw angular function is not itself an $H^{1/2}$ trace.

The area $A_\alpha(h)=|P_\alpha(h)|$ is a quadratic polynomial in $h$.
Indeed, by \eqref{JP:eq:vertex-linear},
\begin{equation}\label{JP:eq:area-jets}
\begin{aligned}
 A_\alpha(h)&=\frac12\sum_i\det(p_i(h),p_{i+1}(h)),\\
 A_\alpha'[h;z]
 &=\frac12\sum_i\{\det(p_i(z),p_{i+1}(h))
                 +\det(p_i(h),p_{i+1}(z))\},\\
 A_\alpha''[z,z]&=\sum_i\det(p_i(z),p_{i+1}(z)).
\end{aligned}
\end{equation}
Equivalently, direct differentiation of the side lengths gives the
ratio-sensitive but exact cyclic formula
\begin{equation}\label{JP:eq:area-second-explicit}
\begin{aligned}
 A_\alpha''[z,z]
 =\sum_i z_i\{&z_{i-1}\csc\alpha_{i-1}
               +z_{i+1}\csc\alpha_i\\
              &-z_i(\cot\alpha_{i-1}+\cot\alpha_i)\}.
\end{aligned}
\end{equation}
Equivalently,
\begin{equation}\label{JP:eq:area-jump}
 A_\alpha''[z,z]=\frac12\sum_i\left\{
 \tan\frac{\alpha_i}{2}(z_i+z_{i+1})^2
 -\cot\frac{\alpha_i}{2}(z_{i+1}-z_i)^2\right\}.
\end{equation}
Thus the area row contains a negative jump channel; it may not be bounded
from below independently of the endpoint part of $\Dred$.

Let $h_0=h_\alpha\mathbf1$ be the canonical tangential support vector,
$A_0=A_\alpha(h_0)$, and define the linear scale functional
\begin{equation}\label{JP:eq:scale-section}
 s_\alpha(h)=\frac{A_\alpha'[h_0;h]}{2A_0}.
\end{equation}
Euler's identity for the quadratic area gives $s_\alpha(h_0)=1$, while
translation invariance of area makes $s_\alpha$ vanish on translation
vectors.  Fix once and for all a translation complement
$\mathsf Y_\alpha$ which contains $h_0$, and put
\begin{equation}\label{JP:eq:Xalpha}
 \mathsf X_\alpha=\{z\in\mathsf Y_\alpha:s_\alpha(z)=0\}.
\end{equation}
Then $h_0+tz$ stays in the affine section $s_\alpha=1$ for every
$z\in\mathsf X_\alpha$.  All support variables, Hessians, dual norms, and
Riesz maps used in JP--SC/JP--JE/JP--FS refer to this fixed section.

To match the global wrapper, let $\mathcal U_{S_0,\rho_0}$ denote the
normalized radial endpoint chart: $S=a+\max_i\alpha_i\le S_0$, all sides
are active, and, after the fixed scale/translation normalization, the
radial graph of $P_\alpha(h)$ is within $\rho_0$ of the unit circle in
$W^{1,\infty}$.  The radii $S_0,\rho_0$ are universal and may be
decreased.  No smallness of the affine-side parametrization relative to
$P_\alpha(h_0)$ is imposed: such a condition would incorrectly exclude a
positive side whose length is much smaller than its normal gap.  The
ratio-free compatibility between the radial endpoint chart, the whole
fixed-support segment, and the support metric is proved below as JP--NS.

\begin{proposition}[The radial chart does not imply a labelled affine-side chart]
\label{JP:prop:radial-not-affine}
There are active area-$\pi$ convex $N$-gons $Q_N$ with positive normal
fans such that
\[
 \max_i\alpha_i+\frac{2\pi}{N}\longrightarrow0,
 \qquad \|r_{Q_N}-1\|_{W^{1,\infty}(\mathbb S^1)}\longrightarrow0,
\]
whereas the derivative of the labelled physical-affine displacement from
the tangential polygon with the same fan tends to infinity.  Thus no
ratio-free radial-to-affine-side chart implication is available.
\end{proposition}

\begin{proof}
Put $a=2\pi/N$, $\varepsilon=a^4$,
$b=(2\pi-2\varepsilon)/(N-2)$, and $t=a^2$.  Choose consecutive normals
$-\varepsilon,0,\varepsilon$ and make all remaining gaps equal to $b$.
For
\[
 c=\left\{\frac{\pi}
 {2\tan(\varepsilon/2)+(N-2)\tan(b/2)}\right\}^{1/2}
 =1+O(a^2),
\]
the constant support vector $c\mathbf1$ gives the area-$\pi$ tangential
polygon.  Keep the central support equal to $c$, change the two adjacent
supports to $c\cos\varepsilon+t\sin\varepsilon$, and leave all other
supports equal to $c$.  The central side then has endpoints $(c,-t)$ and
$(c,t)$ and length $2t$, whereas its tangential reference length is
$2c\tan(\varepsilon/2)$.

Writing
\[
 \delta=\sin\varepsilon
 \left(t-c\tan\frac\varepsilon2\right)
 =t\varepsilon(1+o(1)),
\]
the two incident side lengths change by
$-\delta(\cot\varepsilon+\cot b)=O(a^2)$ from positive lengths
asymptotic to $a$, the next two change by $\delta\csc b$, and all others
are unchanged.  Hence all sides remain active.  The affine scale-section
factor is $1+O(bt\varepsilon)$, so scale and area normalization do not
alter the following asymptotics.  On the central labelled side the
tangential derivative of the displacement is at least
\[
 \frac{t}{(1+o(1))c\tan(\varepsilon/2)}-1
 \sim\frac{2t}{\varepsilon}=\frac2{a^2}\longrightarrow\infty.
\]

Every changed vertex is nevertheless displaced by $O(t)$; the two
central displacements are exactly
$t-c\tan(\varepsilon/2)$, and the outer displacements are
$O(\delta/\sin b)$.  On a side with normal $\theta_i$ the radial graph is
$r(\theta)=h_i/\cos(\theta-\theta_i)$ and
$|r'|=r|\tan(\theta-\theta_i)|$.  Thus, after the harmless area
normalization,
$\|r_{Q_N}-1\|_{W^{1,\infty}}\le Ca\to0$.  This proves the claim.
\end{proof}

\subsection{Exact eigenvalue and scale-normalized Hessians}

Let $u_t>0$ be the $L^2$-normalized first eigenfunction of
$P_\alpha(h+tz)$, let $\lambda_t$ be its eigenvalue, and put
$p_t=-\partial_{\nu_t}u_t$ on the open sides.  If $g$ is boundary data,
write $\Dred_{P_t}(g)$ for the reduced Helmholtz energy: solve the
Eulerian state equation with trace $-g$, impose $L^2$ orthogonality to
$u_t$, and evaluate
\(
 \int_{P_t}|\nabla\psi|^2-\lambda_t\int_{P_t}\psi^2.
\)
This is the same reduced form as in the polygon-preserving volume Hessian.

\begin{theorem}[JP--EH: exact fixed-normal support Hessian]
\label{JP:thm:EH}
For every active support line $h(t)=h+tz$,
\begin{align}
 \lambda_t'&=-\sum_i z_i\int_{S_i(t)}p_t^2\,ds,
 \label{JP:eq:lambda-first}\\
 \lambda_t''&=2\Dred_{P_t}(z_i\partial_{\nu_t}u_t).
 \label{JP:eq:lambda-second}
\end{align}
Hence, with $A_t=A_\alpha(h(t))$ and
\(
 F_t(z)=\sum_i z_i\int_{S_i(t)}p_t^2\,ds,
\)
one has the exact identity
\begin{equation}\label{JP:eq:Phi-second}
 \boxed{
 \frac{d^2}{dt^2}\Phi_\alpha(h(t))
 =\frac1\pi\left\{
 2A_t\Dred_{P_t}(z_i\partial_{\nu_t}u_t)
 -2A_t'[z]F_t(z)+\lambda_t A_\alpha''[z,z]
 \right\}.}
\end{equation}
No stationarity, minimum side, adjacent-ratio, moving radial knot, or
Taylor truncation is used in this formula.
\end{theorem}

\begin{proof}
Hadamard's formula and \eqref{JP:eq:constant-normal} give
\eqref{JP:eq:lambda-first}.  Differentiate it on the moving polygon.  If
$\psi_t$ is the Eulerian eigenfunction derivative, the material derivative
of the boundary flux is
\[
 \dot p_t=-\partial_{\nu_t}\psi_t+V_\tau\partial_\tau p_t.
\]
The differentiated arclength, normal, and transport rows combine, after
one sidewise integration by parts, into
\[
 \lambda_t''=2\Dred_{P_t}((V\cdot\nu_t)\partial_{\nu_t}u_t)
 +2\sum_i\int_{S_i(t)}p_t^2V_\tau
          \partial_\tau(V\cdot\nu_i)\,ds.
\]
The endpoint terms vanish by the convex-corner expansion of the ground
state.  The second row is zero by \eqref{JP:eq:constant-normal}, proving
\eqref{JP:eq:lambda-second}.  Finally differentiate
$\Phi_\alpha=A_t\lambda_t/\pi$ twice, use
$\lambda_t'=-F_t(z)$ and \eqref{JP:eq:area-jets}, and obtain
\eqref{JP:eq:Phi-second}.
\end{proof}

\subsection{Why the disk cubic truncation cannot be the global metric}

For reference, let
\begin{equation}\label{JP:eq:disk-cubic-candidate}
 H^{(3)}_\alpha(z,z)
 =D^2\mathcal L(0)[Z_\alpha,Z_\alpha]
  +D^3\mathcal L(0)[B_\alpha,Z_\alpha,Z_\alpha]
\end{equation}
be the formerly proposed disk-based support Hessian.  The complete physical
chain rule is understood in \eqref{JP:eq:disk-cubic-candidate}; in particular
the two curvature terms are not omitted.

\begin{proposition}[Alternating-mesh obstruction to the cubic metric]
\label{JP:prop:cubic-counterexample}
There are even cyclic fans with fixed adjacent ratio and admissible
translation/area-normalized support directions for which
$H^{(3)}_\alpha(z,z)<0$.  Consequently no positive no-ratio metric can be
obtained by proving coercivity of \eqref{JP:eq:disk-cubic-candidate}.
\end{proposition}

\begin{proof}
Let $N=2M$, $a=2\pi/N$, take alternating exterior angles
$a/10,19a/10$, and set $z_{2j}=1$, $z_{2j+1}=-1$.  The alternating
symmetry gives the two translation constraints and the exact area row.
After the material rescaling $y=\theta/a$, one period has length $2$.
Put $p=(1/10,19/10)$, $c_0=0$, $c_1=1/10$, and on the cell
$x=y-c_r\in[-p_{r-1}/2,p_r/2]$ define
\[
 F_r=\frac{x^2}{2}-\frac{p_0^3+p_1^3}{48},\qquad
 W_0=-\frac{360}{19}-\frac{400}{19}x,\qquad
 W_1=-\frac{360}{19}+\frac{400}{19}x,
\]
\[
 V_r=z_r+xW_r,\qquad U_r=xW_r,\qquad
 K_r=W_r^2-2W_rV_r'.
\]
For period-two Fourier coefficients
$\widehat h(n)=\frac12\int_0^2h(y)e^{-\pi iny}\,dy$, write
\(
 \langle g,h\rangle_2
 =\sum_{n\ne0}\pi|n|\overline{\widehat g(n)}\widehat h(n).
\)
The small-cell expansions, in the physical angular coordinate, are
\[
 f=a^2F+O(a^4),\quad w=a^{-1}W+O(a),\quad
 v=V+O(a^2),\quad f'=aF'+O(a^3),
\]
\[
 w^2-2wv'=a^{-2}K+O(1),\qquad wf'=U+O(a^2).
\]
The exact physical-chain identity is
\[
 D^3\mathcal L(0)[B,Z,Z]
 =D^3\Lambda^n[f,v,v]
  +D^2\Lambda^n[f,w^2-2wv']
  +2D^2\Lambda^n[v,-wf'],
\]
and $D^2\Lambda^n=2q_{\mathbb D}$.  The normal cubic is $O(1)$;
rescaling the other three pairings therefore gives
\[
 H^{(3)}_\alpha(z,z)=\frac{4\pi d^2}{a}E+O(1),\qquad
 E=\langle V,V\rangle_2+\langle F,K\rangle_2
       -2\langle V,U\rangle_2,
\]
where $d>0$ is the disk normalization constant.  Exact rational-polynomial
Fourier integration through mode $4096$, with outward intervals of radius
$10^{-12}$ including the trigonometric Taylor and rounding remainders,
gives
\begin{align*}
 A_{4096}&=46.5557803349116641,\\
 B_{4096}&=22.9387363887007744,\\
 C_{4096}&=35.1293696664596509,
\end{align*}
and hence $E_{4096}=A_{4096}+B_{4096}-2C_{4096}
=-0.7642226093068633$.  Integration by parts in each Fourier coefficient
and Cauchy--Schwarz give the explicit tail bound
\[
 |E-E_K|\le
 \frac{\operatorname{TV}(V')^2}{4\pi^3K^2}
 +\frac{\operatorname{TV}(F')\operatorname{TV}(K)}{2\pi^2K}
 +\frac{\operatorname{TV}(V')\operatorname{TV}(U)}{\pi^2K}.
\]
Here
\begin{align*}
 \operatorname{TV}(F')&=4,&
 \operatorname{TV}(V')&=1600/19,\\
 \operatorname{TV}(U)&=476/19,&
 \operatorname{TV}(K)&=523200/361.
\end{align*}
The reproducible certificate is
\nolinkurl{tests/check_jp_cubic_counterexample.py}.  At $K=4096$ the
right-hand side is less than $0.123892271$, so
$E<-0.640330338$.  Therefore \eqref{JP:eq:disk-cubic-candidate} is negative
for all sufficiently small $a$.  Reflection through a common cell vertex
makes $B$ even and $Z$ odd, so the old linear forcing in this direction and
$D^3\mathcal L(0)[Z,Z,Z]$ vanish.  The negative sign is therefore intrinsic
to the candidate Hessian.  This is an obstruction to the truncation, not
to the exact Hessian \eqref{JP:eq:Phi-second}.
\end{proof}

\begin{corollary}[Stationary specialization]
\label{JP:cor:stationary}
If $h$ is stationary for $\Phi_\alpha$ under the fixed-normal support
variations, then at $t=0$
\begin{equation}\label{JP:eq:stationary-specialization}
 D^2\Phi_\alpha(h)[z,z]
 =\frac{2A}{\pi}\Dred_P(z_i\partial_\nu u)
 +\frac\lambda\pi\left\{A_\alpha''[z,z]
       -\frac{2A_\alpha'[z]^2}{A}\right\}.
\end{equation}
On the area tangent row $A_\alpha'[z]=0$, the last square vanishes.
\end{corollary}

\begin{proof}
Stationarity and \eqref{JP:eq:lambda-first} give
$F_0(z)=\lambda A_\alpha'[z]/A$.  Substitute this in
\eqref{JP:eq:Phi-second}.
\end{proof}

\subsection{Exact value integration and the reduced frontier}

Let $h_\alpha\mathbf1$ denote the canonical tangential support vector and
let $\alpha_*=a\mathbf1$, $a=2\pi/N$.  Put
\begin{equation}\label{JP:eq:forcing}
 \ell_\alpha=D_h\Phi_\alpha(h_\alpha\mathbf1)
\end{equation}
restricted to $\mathsf X_\alpha$.  Ordinary one-variable Taylor
integration, using the genuine Hessian rather than its disk cubic
truncation, gives the identity
\begin{equation}\label{JP:eq:exact-integration}
 \boxed{
 \Phi_\alpha(h_\alpha\mathbf1+z)-\Phi_\alpha(h_\alpha\mathbf1)
 =\ell_\alpha[z]
 +\int_0^1(1-t)
 D_h^2\Phi_\alpha(h_\alpha\mathbf1+tz)[z,z]\,dt.}
\end{equation}

The exact Hadamard variable is the constant side-normal speed $z_i$, not
the radial component of the full material vertex field.  Let
$\Pi_{\rm ng}$ delete the constant and first Fourier harmonics.  We import
from the CT core the literal disk Bessel form
\[
 q_D(f)=d^2\int_{\mathbb T}(\Pi_{\rm ng}f)
       (\mathcal D_\circ+1)(\Pi_{\rm ng}f),
 \qquad
 q_D(f)\asymp\|\Pi_{\rm ng}f\|_{H^{1/2}}^2,
\]
where $d=\partial_ru_0(1)$ and
$\mathcal D_\circ e^{in\theta}
=jJ_{|n|}'(j)J_{|n|}(j)^{-1}e^{in\theta}$ for $n\ne0$.
Let $\theta_{i+1}-\theta_i=\alpha_i$ and define the periodic fan-normal sine
interpolant by
\begin{equation}\label{JP:eq:normal-speed-interpolant}
 (\mathcal I_\alpha z)(\theta)
 =\frac{\sin(\theta_{i+1}-\theta)z_i
         +\sin(\theta-\theta_i)z_{i+1}}{\sin\alpha_i},\qquad
 \theta_i\le\theta\le\theta_{i+1}.
\end{equation}
Put
\begin{equation}\label{JP:eq:normal-speed-metric}
 \mathfrak G_\alpha^{\rm ns}(z)
 :=q_D(\Pi_{\rm ng}\mathcal I_\alpha z),
 \qquad z\in\mathsf X_\alpha.
\end{equation}
This is the metric interface \(\mathrm{JP\text{-}NS}\).

\begin{proposition}[Closed chart entry for the normal-speed metric]
\label{JP:prop:normal-speed-chart-entry}
The form $\mathfrak G_\alpha^{\rm ns}$ is a Hilbert form on
$\mathsf X_\alpha$.  There is a universal $C$ such that
\begin{equation}\label{JP:eq:normal-speed-chart-entry}
 P_\alpha(h_0+z)\in\mathcal U_{S_0,\rho_0}
 \quad\Longrightarrow\quad
 h_0+[0,1]z\ \hbox{is active},\qquad
 \mathfrak G_\alpha^{\rm ns}(z)^{1/2}\le C(S_0+\rho_0),
\end{equation}
and every polygon on the segment has a radial graph $1+\rho_t$ satisfying
\begin{equation}\label{JP:eq:radial-path-entry}
 \sup_{0\le t\le1}\|\rho_t\|_{W^{1,\infty}(\mathbb S^1)}
 \le C(S_0+\rho_0).
\end{equation}
No minimum angle or adjacent-ratio bound is used.
\end{proposition}

\begin{proof}
If $\mathfrak G_\alpha^{\rm ns}(z)=0$, then
$\mathcal I_\alpha z$ contains only the constant and first Fourier modes.
On every open normal-angle interval it satisfies
$(\mathcal I_\alpha z)''+\mathcal I_\alpha z=0$.  Applying this equation
to a function of the form $c+b\cdot e_r(\theta)$ forces $c=0$.  Hence
$z_i=b\cdot\nu_i$ is an exact support translation, which is zero in the
chosen complement $\mathsf X_\alpha$.  This proves positive definiteness.
It also shows that \eqref{JP:eq:normal-speed-metric} descends exactly through
the translation quotient, rather than merely after choosing the complement.

Let $P_i(z)$ be the displacement of the vertex shared by sides $i$ and
$i+1$ from its position in $P_\alpha(h_0)$.  If a star-shaped boundary is
$r(\varphi)e_r(\varphi)$, its one-sided outer normal is parallel to
$r e_r-r'e_\varphi$.  Hence $\|r-1\|_{W^{1,\infty}}\le\rho_0<1/2$
puts the polar direction of each endpoint vertex within $C\rho_0$ of
its adjacent normal cone.  The tangential reference has
$\|r_0-1\|_{W^{1,\infty}}\le CS_0$, and the cone width is at most $S_0$.
The two vertex radii are bounded above and below universally.  Therefore
\begin{equation}\label{JP:eq:vertex-displacement-radial}
 \max_i|P_i(z)|\le C(S_0+\rho_0),
\end{equation}
with no inverse angle or side length.  At that vertex the
constant-normal-speed identity gives
\[
 z_i=P_i(z)\cdot\nu_i,\qquad
 z_{i+1}=P_i(z)\cdot\nu_{i+1}.
\]
Solving these two projection equations gives the exact identity
\[
 (\mathcal I_\alpha z)(\theta)=P_i(z)\cdot e_r(\theta),
 \qquad \theta_i\le\theta\le\theta_{i+1}.
\]
Therefore
\[
 \int_{\theta_i}^{\theta_{i+1}}
 \left\{|\mathcal I_\alpha z|^2
       +|\partial_\theta\mathcal I_\alpha z|^2\right\}\,d\theta
 =\alpha_i|P_i(z)|^2.
\]
Summation gives
$\|\mathcal I_\alpha z\|_{H^1}\le C(S_0+\rho_0)$, and the
Bessel equivalence $q_D\asymp H^{1/2}$ proves the metric bound.  For fixed
normals every side length is a linear function of the support numbers.
It is positive at both $h_0$ and $h_0+z$; therefore it remains positive on
their whole affine segment.  This proves activity of every labelled side.

Every intermediate vertex is the affine combination of its endpoint
vertices.  Its radial and tangential coordinates relative to either
incident normal therefore obey the same $C(S_0+\rho_0)$ bounds.  On each
side the exact formulas
\[
 r_t(\theta)=\frac{h_i(t)}{\cos(\theta-\theta_i)},
 \qquad
 r_t'(\theta)=r_t(\theta)\tan(\theta-\theta_i)
\]
then give~\eqref{JP:eq:radial-path-entry}, uniformly in $t$.  In view of
Proposition~\ref{JP:prop:radial-not-affine}, this radial/metric path statement,
not a labelled affine-side derivative bound, is the ratio-free interface
used downstream.
\end{proof}

\begin{proposition}[Exact area identity in the normal-speed chart]
\label{JP:prop:normal-speed-area}
For $f=\mathcal I_\alpha z$ one has
\begin{equation}\label{JP:eq:normal-speed-area}
 \boxed{A_\alpha''[z,z]
 =\int_{\mathbb T}\bigl(f^2-|f'|^2\bigr)\,d\theta.}
\end{equation}
Moreover, on each interval $I_i=[\theta_i,\theta_{i+1}]$,
\begin{equation}\label{JP:eq:normal-speed-L2}
\begin{aligned}
 2\int_{I_i}f^2\,d\theta
 ={}&\frac{\alpha_i-\sin\alpha_i}{4\sin^2(\alpha_i/2)}
       (z_{i+1}-z_i)^2\\
 &+\frac{\alpha_i+\sin\alpha_i}{4\cos^2(\alpha_i/2)}
       (z_{i+1}+z_i)^2.
\end{aligned}
\end{equation}
Thus the two displayed coefficients are exactly $2\int_{I_i}f^2$.  After
the Hessian factor $\lambda_\alpha/\pi$, this is the polygonal local
$L^2$ row converging to the $L^2$ part of $2q_D$.  The low-mode projection
and the nonlocal order-one row of $q_D$ remain global.
\end{proposition}

\begin{proof}
Write $s=\theta-(\theta_i+\theta_{i+1})/2$ and
\[
 X_i=\frac{z_i+z_{i+1}}{2\cos(\alpha_i/2)},
 \qquad
 Y_i=\frac{z_{i+1}-z_i}{2\sin(\alpha_i/2)}.
\]
The sine interpolation formula is exactly
$f=X_i\cos s+Y_i\sin s$.  Orthogonality of the even and odd summands on
$[-\alpha_i/2,\alpha_i/2]$ gives \eqref{JP:eq:normal-speed-L2}.  The same
calculation gives
\[
 \int_{I_i}(f^2-|f'|^2)\,d\theta
 =2z_iz_{i+1}\csc\alpha_i
  -(z_i^2+z_{i+1}^2)\cot\alpha_i.
\]
Summing in $i$ is precisely \eqref{JP:eq:area-second-explicit}, which proves
\eqref{JP:eq:normal-speed-area}.
\end{proof}

Henceforth the support Hilbert form used by JP--SC, JP--FS, and NR is fixed
to be
\begin{equation}\label{JP:eq:JP-metric-fixed}
 \Gjp_\alpha:=\mathfrak G_\alpha^{\rm ns}.
\end{equation}

\begin{remark}[Archived stronger JP--SC route]
\label{JP:rem:SC}
For every fixed $\delta\in(0,1)$ there are
$S_0,\rho_0,\eta_0>0$, independent of the number, minimum size, and ratios
of the sides, such that the following holds for every fan with
$S=a+\max_i\alpha_i\le S_0$.  For all
$y,\xi\in\mathsf X_\alpha$ with the
entire segment $h_0+[0,1]y$ in the physical active support chart and
$\mathfrak G_\alpha^{\rm ns}(y)^{1/2}\le\eta_0$,
\begin{equation}\label{JP:eq:SC}
 D_h^2\Phi_\alpha(h_\alpha\mathbf1+y)[\xi,\xi]
 \ge (2-\delta)\mathfrak G_\alpha^{\rm ns}(\xi).
\end{equation}
The exported Hilbert constant is $c=2-\delta$; it may not be silently
decreased before JP--FS.
The proof must use the complete expression \eqref{JP:eq:Phi-second}.  A lower
bound for its reduced-energy diagonal alone is insufficient unless the
area and mixed-flux rows have first been Schur-completed.
\end{remark}

\begin{remark}[Why stationary O--H does not discharge JP--SC]
At the base point $h_\alpha\mathbf1$, the definition of the fixed scale
section gives $A_\alpha'[\xi]=0$, so the exact base Hessian is
\begin{equation}\label{JP:eq:base-support-hessian}
 D_h^2\Phi_\alpha(h_\alpha\mathbf1)[\xi,\xi]
 =\frac1\pi\left\{
 2A_\alpha\Dred_{P_\alpha^0}(\xi_i\partial_\nu u_\alpha)
 +\lambda_\alpha A_\alpha''[\xi,\xi]\right\}.
\end{equation}
The stationary O--H comparison is neither stated at this generally
nonstationary tangential polygon nor relative to the scalar support metric.
Its available error has the form
$\omega(S)\|V\|_{H^{1/2}(\mathbb T;\mathbb R^2)}^2$.  This full-vector
norm is not uniformly controlled by \eqref{JP:eq:normal-speed-metric}: the
vertex solution contains the tangential coefficient
\[
 \frac{\xi_{i+1}-\xi_i\cos\alpha_i}{\sin\alpha_i},
\]
which has an inverse-angle amplification that the scale-critical
fan-normal interpolation does not.  A rate-free $\omega(S)\to0$ therefore
cannot be absorbed into $\mathfrak G_\alpha^{\rm ns}$.  A proof of JP--SC
needs a support-compatible,
endpoint-resummed relative lower bound for the complete form
\eqref{JP:eq:base-support-hessian}, followed by control of the mixed-flux row
away from the base.  Invoking stationary O--H or its full-vector norm is
not such a bound.
\end{remark}

\begin{remark}[Why a pointwise corner-amplitude estimate is insufficient]
The terminal jump--cotangent cancellation in the stationary O--H ledger
uses sidewise stationary flux normalization.  At the generally
nonstationary tangential polygon one has only the global Rellich average,
not
\(
 \int_{S_i}p^2=(\lambda/\pi)|S_i|
\)
for each side.  If $\beta_i$ denotes the normalized leading flux amplitude
at a corner of exterior angle $\alpha_i$, separating that corner before
its descendants leaves a schematic residual
\[
 (\beta_i^2-1)\alpha_i^{-1}(z_{i+1}-z_i)^2.
\]
The corresponding sharp transition in the JP--NS form costs only
$O((z_{i+1}-z_i)^2(1+|\log\alpha_i|))$.  Consequently even a uniform
$\beta_i=1+o(1)$ is not an absorbable pointwise error.  The amplitude loss
must be retained with all descendant annuli and LCA pairs until their
positive energy is recovered.  This is why \eqref{JP:eq:graded-DtN-blocker}
is a completed operator inequality rather than a cornerwise estimate.
\end{remark}

The fan-normal candidate reduces JP--SC to two precise estimates.  In the
pure two-ray wedge model, combining the endpoint reduced energy with the
area row gives exactly the local integral identity
\eqref{JP:eq:normal-speed-L2}.
Both coefficients are nonnegative, and the common coefficient is positive.
What is not supplied by this local calculation is the ratio-free global
normalization of all terminal rays, annuli, and LCA pairs.

\begin{remark}[Archived JP--S0 sufficient estimate]
\label{JP:rem:S0}
Prove a modulus $\omega_0(S)\to0$, independent of all minimum angles and
adjacent ratios, such that
\begin{equation}\label{JP:eq:S0}
 \boxed{
 D_h^2\Phi_\alpha(h_0)[z,z]
 \ge(2-\omega_0(S))\mathfrak G_\alpha^{\rm ns}(z),
 \qquad z\in\mathsf X_\alpha.}
\end{equation}
The complete two-ray, annular, and LCA ledger must be retained, and its
amplified jump must be combined with $\lambda_0A_\alpha''[z,z]$ before
either term is estimated.
\end{remark}

At the base scale section, \eqref{JP:eq:normal-speed-area} rewrites this
obligation without any hidden geometric row as
\begin{equation}\label{JP:eq:graded-DtN-blocker}
 \begin{aligned}
 \frac1\pi\left\{
 2A_\alpha\Dred_{P_\alpha^0}(z_i\partial_\nu u_\alpha)
 +\lambda_\alpha\int_{\mathbb T}(f^2-|f'|^2)\right\}
 &\ge(2-\omega_0(S))q_D(\Pi_{\rm ng}f),\\
 f&=\mathcal I_\alpha z.
 \end{aligned}
\end{equation}
The same/adjacent two-ray part is \eqref{JP:eq:normal-speed-L2}; the genuinely
missing part is the no-ratio nonlocal terminal-ray, annular, and LCA
normalization in \eqref{JP:eq:graded-DtN-blocker}.  None of H0, QI, DF, or DM
states this polygonal reduced-DtN operator inequality.

\begin{remark}[Archived JP--ST sufficient estimate]
\label{JP:rem:ST}
Along every active fixed-scale support segment in the physical chart,
prove a modulus $\omega_1(\rho,S)\to0$ such that
\begin{equation}\label{JP:eq:ST}
 D_h^2\Phi_\alpha(h_0+y)[\xi,\xi]
 \ge\min\left\{D_h^2\Phi_\alpha(h_0)[\xi,\xi],
                 2\mathfrak G_\alpha^{\rm ns}(\xi)\right\}
       -\omega_1(\rho,S)\mathfrak G_\alpha^{\rm ns}(\xi).
\end{equation}
All reduced-energy, mixed-flux, and area rows in
\eqref{JP:eq:Phi-second} are included.  No two-sided absolute Hessian
transport is required: a large positive endpoint coefficient need only
remain above the disk threshold, not be approximated to additive
$\mathfrak G_\alpha^{\rm ns}$ accuracy.
\end{remark}

\begin{proposition}[Typed sufficient cut for JP--SC]
\label{JP:prop:SC-cut}
JP--S0 and JP--ST, together with
Proposition~\ref{JP:prop:normal-speed-chart-entry}, imply JP--SC with
$\Gjp_\alpha=\mathfrak G_\alpha^{\rm ns}$ and
$c=2-\delta$ after decreasing $S_0$ and $\rho_0$ so that
$\omega_0(S)+\omega_1(\rho,S)\le\delta$.
\end{proposition}

\begin{proof}
Choose the thresholds so that
$\omega_0(S)+\omega_1(\rho,S)\le\delta$.  Then
\eqref{JP:eq:S0} makes the minimum in \eqref{JP:eq:ST} at least
$(2-\omega_0(S))\mathfrak G_\alpha^{\rm ns}(\xi)$.  Hence
\eqref{JP:eq:S0}--\eqref{JP:eq:ST} give
$D_h^2\Phi_\alpha(h_0+y)[\xi,\xi]
\ge(2-\delta)\mathfrak G_\alpha^{\rm ns}(\xi)$.  The remaining chart and activity
requirements are exactly \eqref{JP:eq:normal-speed-chart-entry}.
\end{proof}

\begin{remark}[Numerical warning is not a counterexample]
Finite-element prechecks on alternating fans have not found a negative
exact support Hessian.  The quotient relative to the older physical-radial
trace deteriorates as one angle ratio decreases at fixed modest $N$; that
calculation may also underresolve the nearly flat corner exponent and is
not a certified asymptotic.  On the same samples, the quotient relative to
$\mathfrak G_\alpha^{\rm ns}$ does not deteriorate.  These are design
diagnostics only: \eqref{JP:eq:S0} still requires a ratio-free proof for the
complete endpoint ledger, and no node status follows from the precheck.
\end{remark}

\begin{remark}[Status of the candidate cut]
JP--S0 and JP--ST are typed sufficient subcontracts; they are not
additional active DAG vertices.  A replacement proof using another
no-ratio support metric would have to replace JP--NS and the JP--SC
interface together.  Conversely, a rigorous degeneration of JP--S0 would
reject only this particular completion route, not the target theorem.
\end{remark}

Let $\|\ell_\alpha\|_{(\Gjp_\alpha)^*}$ be the dual norm.  We first remove
the part of the former JP--BR contract which is already a consequence of
the closed pure-fan nodes.

\begin{theorem}[JP--BU: exact tangential-fan baseline upper bound]
\label{JP:thm:BU}
There are universal $C,S_0>0$ such that, for every positive fan with
$S=a+\max_i\alpha_i\le S_0$,
\begin{equation}\label{JP:eq:fan-upper}
 \boxed{
 \Phi_\alpha(h_\alpha\mathbf1)
       -\Phi_{\alpha_*}(h_{\alpha_*}\mathbf1)
 \le C\Jfour(\alpha).}
\end{equation}
The constant is independent of the smallest angle and all adjacent ratios.
\end{theorem}

\begin{proof}
The exact disk-normalized expansion furnished by the closed TR, CT core,
QB, and DF nodes splits the pure tangential-fan value into its quadratic
disk-source row, its complete cubic row, and a fourth-and-higher scalar
tail $F_4$; explicitly, if $\Delta_2$ and $\Delta_3$ denote the complete
degree-two and degree-three scalar rows, then
\begin{align*}
 \Phi_\alpha(h_\alpha\mathbf1)
 -\Phi_{\alpha_*}(h_{\alpha_*}\mathbf1)
 &=[\Delta_2(\alpha)-\Delta_2(\alpha_*)]
   +[\Delta_3(\alpha)-\Delta_3(\alpha_*)]\\
 &\quad+[F_4(\alpha)-F_4(\alpha_*)].
\end{align*}
The relative source transfer gives the required one-sided quadratic
Bregman estimate, while the closed pure cubic row is absolute:
\begin{align*}
 \Delta_2(\alpha)-\Delta_2(\alpha_*)&\le C\Jfour(\alpha),\\
 |\Delta_3(\alpha)-\Delta_3(\alpha_*)|&\le C\Jfour(\alpha).
\end{align*}
It remains only to record why the scalar tail needs no support-direction
CT--M4 estimate.

Put $d_i=\alpha_i-a$ and
$\gamma_i(t)=a+td_i$.  The unconditional tame scalar fourth-remainder
corollary of the closed CT core (obtained by applying Cauchy's coefficient
estimate to the complex natural-order Hessian) gives, on the real chart
and on its material complex polydisc,
\begin{equation}\label{JP:eq:F4-absolute}
 |F_4(\zeta)|
 \le C\|B_\zeta\|_{H^{1/2}}^2
        \|B_\zeta\|_{W^{1,\infty}}^2
 \le CS^5\le CS^4.
\end{equation}
Choose a fixed sufficiently small $\varepsilon>0$.  If
$\Jfour(\alpha)\ge\varepsilon S^4$, applying
\eqref{JP:eq:F4-absolute} at the two endpoints gives the desired bound with
constant $2C/\varepsilon$.  If
$\Jfour(\alpha)<\varepsilon S^4$, the exact identity
\begin{equation}\label{JP:eq:J4-action}
 \Jfour(\alpha)
 =12\int_0^1(1-t)E(t)\,dt,
 \qquad E(t)=\sum_i\gamma_i(t)^2d_i^2,
\end{equation}
and
\(
 \Jfour=\sum_i d_i^2(\alpha_i^2+2a\alpha_i+3a^2)
\)
show, after decreasing $\varepsilon$, that
$a\asymp S$ and the whole segment is uniformly comparable.  The material
complex radius in the direction $d$ is at least
$ca/\|d\|_{\ell^\infty}$.  Cauchy's estimate and
\eqref{JP:eq:F4-absolute} therefore give
\[
 |F_4''(t)|
 \le CS^4\frac{\|d\|_{\ell^\infty}^2}{a^2}
 \le C E(t).
\]
Cyclic symmetry makes $F_4'(0)$ constant in the fan coordinates, hence
$F_4'(0)[d]=0$ because $\sum_i d_i=0$.  Taylor integration and
\eqref{JP:eq:J4-action} yield
\[
 |F_4(\alpha)-F_4(\alpha_*)|
 \le C\int_0^1(1-t)E(t)\,dt
 \le C\Jfour(\alpha).
\]
Adding the three rows proves \eqref{JP:eq:fan-upper}.  Only the pure scalar
tail was used here; none of the retired support-linear CT--M4 estimates is
being reintroduced.
\end{proof}

The remaining forcing admits a coordinate-free finite-dimensional
description.  Let
\begin{equation}\label{JP:eq:side-masses}
 P_\alpha^0=P_\alpha(h_\alpha\mathbf1),\qquad
 L_i=|S_i|,\qquad
 m_i=\int_{S_i}(-\partial_\nu u_\alpha)^2\,ds,
\end{equation}
where $|P_\alpha^0|=\pi$ and $u_\alpha$ is $L^2$-normalized.

\begin{proposition}[Exact flux representation of the support forcing]
\label{JP:prop:flux-forcing}
On the fixed affine scale section, equivalently
\(
 \sum_iL_i z_i=0
\)
at $h_\alpha\mathbf1$, one has
\begin{equation}\label{JP:eq:exact-flux-forcing}
 \boxed{
 \ell_\alpha[z]
 =-\sum_i m_i z_i
 =-\sum_i\left(m_i-\frac{\lambda_\alpha}{\pi}L_i\right)z_i.}
\end{equation}
Moreover, with
\begin{equation}\label{JP:eq:centered-flux}
 r_{\alpha,i}=m_i-\frac{\lambda_\alpha}{\pi}L_i,
\end{equation}
one has
\begin{equation}\label{JP:eq:flux-moments}
 \sum_i r_{\alpha,i}=0,
 \qquad
 \sum_i r_{\alpha,i}\nu_i=0.
\end{equation}
Thus $r_\alpha$ is an exact covector on the scale/translation quotient.
\end{proposition}

\begin{proof}
Hadamard's formula and the first area variation give
\[
 D_h\Phi_\alpha(h_\alpha\mathbf1)[z]
 =\frac1\pi\left\{
 \lambda_\alpha\sum_iL_i z_i-\pi\sum_i m_i z_i\right\}.
\]
The first term vanishes on the fixed scale section.  Rellich's identity
on the tangential polygon and tangentiality of its boundary give
\[
 h_\alpha\sum_i m_i=2\lambda_\alpha,
 \qquad
 \sum_iL_i=\frac{2\pi}{h_\alpha}.
\]
These identities prove the first equation in \eqref{JP:eq:flux-moments} and
allow the centered form in \eqref{JP:eq:exact-flux-forcing}.  Translation
invariance in Hadamard's formula gives
$\sum_i m_i\nu_i=0$, while the boundary normal-balance identity is
$\sum_iL_i\nu_i=0$.  Their difference proves the vector equation in
\eqref{JP:eq:flux-moments}.
\end{proof}

We now isolate the analytic input which is actually missing from the old
TR/DM route.  On the side-normal cell
$K_i=[\theta_i-\alpha_{i-1}/2,\theta_i+\alpha_i/2]$, put
\begin{equation}\label{JP:eq:disk-source-B}
 B_\alpha=\Pi_{\rm ng}\{h_\alpha\sec(\theta-\theta_i)-1\},
 \qquad K_\alpha z=\Pi_{\rm ng}\mathcal I_\alpha z.
\end{equation}
Let $\langle\cdot,\cdot\rangle_D$ be the bilinear disk Bessel form and
define
\begin{align}
 p_\alpha[z]&=2\langle B_\alpha,K_\alpha z\rangle_D,
 &\Gjp_\alpha(z)&=q_D(K_\alpha z),\label{JP:eq:JE-forms}\\
 P_D^{\rm ns}(\alpha)&=
 \inf_{z\in\mathsf X_\alpha}q_D(B_\alpha+K_\alpha z),
 &\Delta_\alpha&=
 \Phi_\alpha(h_\alpha\mathbf1)
 -\Phi_{\alpha_*}(h_{\alpha_*}\mathbf1).
 \label{JP:eq:D0ns-definition}
\end{align}
Cyclic symmetry gives $p_{\alpha_*}=0$ on the regular scale/translation
quotient.

We first close the disk part of the endpoint.  This is not a transfer from
the old physical-radial space: the normal-cell partition admits its own
exact half-strip square.

For $I_i=[\theta_i,\theta_{i+1}]$, let
\begin{equation}\label{JP:eq:normal-linear-spline}
 (\mathcal I_\alpha^0z)(\theta_i+y)
 =z_i+\frac{y}{\alpha_i}(z_{i+1}-z_i),
 \qquad 0\le y\le\alpha_i,
\end{equation}
and let $G_\alpha^{\rm e}$ be the edge-quadratic source
\begin{equation}\label{JP:eq:edge-source-on-normal-cell}
 G_\alpha^{\rm e}(\theta_i+y)
 =\frac12\min\{y,\alpha_i-y\}^2.
\end{equation}
This is precisely $x^2/2$ on the two side cells incident to the polygonal
vertex direction $\theta_i+\alpha_i/2$.  Put
\begin{equation}\label{JP:eq:q-plus-normal}
 q_+(f)=\sum_{n\ne0}(|n|+1)|\widehat f_n|^2,
 \qquad
 L_\alpha^0z=\frac12\sum_i(\alpha_{i-1}+\alpha_i)z_i,
\end{equation}
and
\begin{equation}\label{JP:eq:P-plus-normal-linear}
 P_{+,0}^{\rm ns}(\alpha)
 =\inf_{L_\alpha^0z=0}q_+(G_\alpha^{\rm e}+\mathcal I_\alpha^0z).
\end{equation}
The missing constant in \eqref{JP:eq:q-plus-normal} is understood as the
global mean minimization.

\begin{lemma}[Exact normal-cell half-strip and variance squares]
\label{JP:lem:normal-cell-square}
Let $\mathcal N_i$ be the least Dirichlet energy in the Neumann half-strip
$I_i\times(0,\infty)$ with prescribed bottom trace.  Then
\begin{equation}\label{JP:eq:normal-cell-strip-square}
 \boxed{
 \mathcal N_i(G_\alpha^{\rm e}+\mathcal I_\alpha^0z)
 =\frac{\zeta(3)}{\pi^3}
 \left\{\frac{\alpha_i^4}{4}
       +7(z_{i+1}-z_i)^2\right\}.}
\end{equation}
Moreover, after allowing an independent local mean $C_i$,
\begin{equation}\label{JP:eq:normal-cell-L2-square}
 \boxed{
 \inf_{C_i\in\mathbb R}\int_{I_i}
 |G_\alpha^{\rm e}+\mathcal I_\alpha^0z-C_i|^2
 =\frac{\alpha_i^5}{720}
  +\frac{\alpha_i}{12}(z_{i+1}-z_i)^2.}
\end{equation}
Consequently
\begin{equation}\label{JP:eq:normal-homogeneous-gap}
 P_{+,0}^{\rm ns}(\alpha)-P_{+,0}^{\rm ns}(\alpha_*)
 \ge \frac{\zeta(3)}{8\pi^4}\Jfour(\alpha).
\end{equation}
The estimate is pointwise in the coefficient vector before imposing the
area and translation rows.
\end{lemma}

\begin{proof}
Write $w=\alpha_i$, $t=y/w$, and
\[
 g(t)=\frac12\min\{t,1-t\}^2,
 \qquad \rho=\frac{z_{i+1}-z_i}{w^2}.
\]
Modulo a constant, the normalized bottom trace is
$w^2h(t)$ with $h(t)=g(t)+\rho t$.  Its distributional second derivative
is $h''=1-\delta_{1/2}$.  Twice integrating the cosine coefficient gives,
for $k\ge1$,
\[
 2\int_0^1h(t)\cos(k\pi t)\,dt
 =\frac{2\{\rho((-1)^k-1)+\cos(k\pi/2)\}}{(k\pi)^2}.
\]
Thus the even modes contain only the source and the odd modes only the
slope.  Since
\[
 \sum_{k\ {\rm even}}k^{-3}=\frac18\zeta(3),
 \qquad
 \sum_{k\ {\rm odd}}k^{-3}=\frac78\zeta(3),
\]
conformal rescaling of the half-strip gives
\eqref{JP:eq:normal-cell-strip-square}.

The function $g$ is symmetric about $1/2$, so it is orthogonal in variance
to $t-1/2$.  Direct integration gives
\[
 \int_0^1g=\frac1{24},\qquad
 \int_0^1g^2=\frac1{320},\qquad
 \operatorname {Var}(g)=\frac1{720},\qquad
 \operatorname {Var}(t)=\frac1{12}.
\]
This proves \eqref{JP:eq:normal-cell-L2-square}.

Restriction of the periodic harmonic extension yields
\[
 2\pi\sum_{n\ne0}|n||\widehat f_n|^2
 \ge\sum_i\mathcal N_i(f).
\]
For the $L^2$ row, independent cell means only decrease the infimum relative
to one global mean.  At the regular fan with $z=0$, reflection at every
normal line makes both restrictions equalities and all local means agree.
Subtracting the regular value, using \eqref{JP:eq:normal-cell-strip-square}--
\eqref{JP:eq:normal-cell-L2-square}, and applying convexity of $x^5$ gives
\[
 P_{+,0}^{\rm ns}(\alpha)-P_{+,0}^{\rm ns}(\alpha_*)
 \ge\frac{\zeta(3)}{8\pi^4}
       \left(\sum_i\alpha_i^4-Na^4\right)
 +\frac1{1440\pi}
       \left(\sum_i\alpha_i^5-Na^5\right),
\]
which proves \eqref{JP:eq:normal-homogeneous-gap}.
\end{proof}

Let $Z_\alpha=\ker L_\alpha$, where
\begin{equation}\label{JP:eq:normal-exact-area-row}
 L_\alpha z=\sum_i\left\{
 \tan\frac{\alpha_{i-1}}2+\tan\frac{\alpha_i}2\right\}z_i,
\end{equation}
and define the exact-sine, edge-source value
\begin{equation}\label{JP:eq:P-plus-normal-sine}
 P_{+,\sin}^{\rm e}(\alpha)
 =\inf_{z\in Z_\alpha}
 q_+(G_\alpha^{\rm e}+\mathcal I_\alpha z).
\end{equation}
Translations may be left in $Z_\alpha$: the sine interpolant sends them
exactly to the first harmonics, so the minimizer cancels those modes.  This
is equivalent to applying $\Pi_{\rm ng}$ and then taking the fixed
translation complement.

\begin{lemma}[Relative linear-to-sine area graph]
\label{JP:lem:normal-sine-graph}
There is a modulus $\omega_{\sin}(S)\to0$, independent of all minimum
angles and adjacent ratios, such that
\begin{equation}\label{JP:eq:normal-sine-relative}
 \left|\{P_{+,\sin}^{\rm e}(\alpha)-P_{+,0}^{\rm ns}(\alpha)\}
 -\{P_{+,\sin}^{\rm e}(\alpha_*)-P_{+,0}^{\rm ns}(\alpha_*)\}\right|
 \le\omega_{\sin}(S)\Jfour(\alpha).
\end{equation}
One may take $\omega_{\sin}(S)=C(S^{1/4}+S^{1/2})$.
\end{lemma}

\begin{proof}
On $I_i$, Taylor's formula for the two endpoint basis functions gives
\begin{align*}
 \|\mathcal I_\alpha z-\mathcal I_\alpha^0z\|_{L^\infty(I_i)}
 &\le C\alpha_i^2(|z_i|+|z_{i+1}|),\\
 \|\partial_\theta(\mathcal I_\alpha z-\mathcal I_\alpha^0z)
        \|_{L^\infty(I_i)}
 &\le C\alpha_i(|z_i|+|z_{i+1}|).
\end{align*}
For every constant $c$, direct integration of a linear function gives the
ratio-free nodal sampling estimate
\[
 \sum_i(\alpha_{i-1}+\alpha_i)|z_i-c|^2
 \le C\|\mathcal I_\alpha^0z-c\|_{L^2}^2.
\]
If $L_\alpha^0z=0$, then
$\int_{\mathbb T}\mathcal I_\alpha^0z=L_\alpha^0z=0$, so take $c=0$.
Summing the displayed jets, applying the sampling estimate, and using
$\|e\|_{H^{1/2}}^2\le C\|e\|_{L^2}\|e\|_{H^1}$ gives
\begin{equation}\label{JP:eq:normal-linear-sine-close}
 \|\mathcal I_\alpha z-\mathcal I_\alpha^0z\|_{H^{1/2}}
 \le CS^{3/2}\|\mathcal I_\alpha^0z\|_{H^{1/2}}.
\end{equation}
Indeed, the squared $L^2$ and $H^1$ errors are bounded respectively by
$CS^4\|\mathcal I_\alpha^0z\|_{L^2}^2$ and
$CS^2\|\mathcal I_\alpha^0z\|_{L^2}^2$.

For $z\in\ker L_\alpha^0$, put
\[
 C_\alpha z=z-\frac{L_\alpha z}{L_\alpha\mathbf1}\mathbf1.
\]
Then $C_\alpha z\in Z_\alpha$, and the analogous formula using
$L_\alpha^0$ is its inverse.  Since
\[
 L_\alpha z-L_\alpha^0z
 =\sum_i\left(\tan\frac{\alpha_i}{2}-\frac{\alpha_i}{2}\right)
       (z_i+z_{i+1}),
\]
the sampling estimate and
$|\tan(u/2)-u/2|\le Cu^3$ show that the scalar correction in $C_\alpha$
has size at most $CS^2\|\mathcal I_\alpha^0z\|_{H^{1/2}}$.
The sine trace of the constant nodal vector has uniformly bounded
$H^{1/2}$ norm.  Combining this fact with
\eqref{JP:eq:normal-linear-sine-close} yields
\begin{equation}\label{JP:eq:normal-sine-graph-close}
 \|\mathcal I_\alpha C_\alpha z-\mathcal I_\alpha^0z\|_{H^{1/2}}
 \le CS^{3/2}\|\mathcal I_\alpha^0z\|_{H^{1/2}},
\end{equation}
with constants independent of the smallest interval and all adjacent
ratios.

The source bound
$q_+(G_\alpha^{\rm e})\le C\sum_i\alpha_i^4\le CS^3$
follows by splitting the Slobodeckij integral into same-cell,
adjacent-cell, and dyadically separated-cell pieces.  The elementary
Hilbert-space fact
\[
 \left|\operatorname {dist}(b,V_1)^2
       -\operatorname {dist}(b,V_0)^2\right|
 \le C\eta\|b\|^2
\]
whenever $V_1$ is the graph of a bijection $J:V_0\to V_1$ with
$\|J-I\|\le\eta<1/2$, applied to
\eqref{JP:eq:normal-sine-graph-close}, gives
\begin{equation}\label{JP:eq:normal-sine-absolute}
 |R_{\sin}(\alpha)|\le CS^{9/2},
 \qquad R_{\sin}=P_{+,\sin}^{\rm e}-P_{+,0}^{\rm ns}.
\end{equation}

It remains to make the comparison relative to $\Jfour$.  In a
quasiuniform material band
$(1-\eta_0)a\le\gamma_i\le(1+\eta_0)a$, pull every interval affinely to
the real fan and complexify
$\gamma(t)=\gamma+t\delta$ for
$|t|<ca/\|\delta\|_{\ell^\infty}$.  The pulled-back Gagliardo kernel, the
$L^2$ row, both interpolation graphs, the area row, and the quadratic
source are holomorphic.  The sampling coercivity above keeps the complex
Gram inverse uniformly bounded after the numerical constant $c$ is
decreased.  Thus the proof of \eqref{JP:eq:normal-sine-absolute} applies to
the complex bilinear Schur formula and gives
$|R_{\sin}(\gamma(t))|\le Ca^{9/2}$.  Cauchy's estimate and
$\sum_i\gamma_i^2\delta_i^2\ge ca^2
\|\delta\|_{\ell^\infty}^2$ give
\[
 |D^2R_{\sin}(\gamma)[\delta,\delta]|
 \le Ca^{1/2}\sum_i\gamma_i^2\delta_i^2,
\]

Put $d_i=\alpha_i-a$.  The exact identities
\[
 \Jfour=\sum_i d_i^2(\alpha_i^2+2a\alpha_i+3a^2)
 =12\int_0^1(1-t)\sum_i(a+td_i)^2d_i^2\,dt
\]
imply $\max_i|d_i|^4\le C\Jfour$.  If
$\Jfour\ge S^{17/4}$, subtract the two bounds
\eqref{JP:eq:normal-sine-absolute}.  Otherwise, $a\asymp S$ and
$\max_i|d_i|/a\le CS^{1/16}$, so the whole radial segment is
quasiuniform.  Cyclic symmetry kills
$DR_{\sin}(a\mathbf1)[d]$ because $\sum_i d_i=0$.  Taylor integration and
the Hessian bound give the $CS^{1/2}\Jfour$ estimate in the near regime;
the far regime gives $CS^{1/4}\Jfour$.  This proves
\eqref{JP:eq:normal-sine-relative}.
\end{proof}

\begin{lemma}[Normal-sine Bessel and secant transfer]
\label{JP:lem:normal-bessel-secant}
With $\widehat P_D^{\rm ns}=(2\pi d^2)^{-1}P_D^{\rm ns}$, there is a
modulus $\omega_D(S)\to0$ such that
\begin{equation}\label{JP:eq:normal-bessel-secant-relative}
 \left|\{\widehat P_D^{\rm ns}(\alpha)-P_{+,\sin}^{\rm e}(\alpha)\}
 -\{\widehat P_D^{\rm ns}(\alpha_*)-P_{+,\sin}^{\rm e}(\alpha_*)\}\right|
 \le\omega_D(S)\Jfour(\alpha).
\end{equation}
\end{lemma}

\begin{proof}
Put
\[
 W_\alpha=\Pi_{\rm ng}\mathcal I_\alpha Z_\alpha,
 \qquad
 \widehat q_D=(2\pi d^2)^{-1}q_D.
\]
Because exact support translations belong to $Z_\alpha$ and are carried by
$\mathcal I_\alpha$ to the first harmonics, minimizing $q_+$ over the
unprojected sine space is the same as minimizing it over $W_\alpha$ after
applying $\Pi_{\rm ng}$.

We first construct the required no-ratio quasi-interpolant.  For
$\psi\in H_{\rm ng}^{3/2}:=\Pi_{\rm ng}H^{3/2}(\mathbb T)$, let
$\phi=P_{\le M}\psi$, $M=\lfloor S^{-1}\rfloor$, sample $phi$ at the
normal angles, and use the linear interpolant
\eqref{JP:eq:normal-linear-spline}.  Fourier truncation gives
\begin{equation}\label{JP:eq:normal-QI-tail}
 \|\psi-\phi\|_{H^{1/2}}\le CS\|\psi\|_{H^{3/2}},
 \qquad
 \|\phi''\|_{L^2}\le CS^{-1/2}\|\psi\|_{H^{3/2}}.
\end{equation}
Taylor's integral formula on each interval, followed by summation, gives
\[
 \|\phi-\mathcal I_\alpha^0\phi|_{\{\theta_i\}}\|_{L^2}
 \le CS^2\|\phi''\|_{L^2},\qquad
 \|\phi-\mathcal I_\alpha^0\phi|_{\{\theta_i\}}\|_{H^1}
 \le CS\|\phi''\|_{L^2}.
\]
The constants are independent of cell ratios because each point belongs to
one normal interval and the estimate is charged to that same interval.
Interpolation between $L^2$ and $H^1$, the tail
\eqref{JP:eq:normal-QI-tail}, the exact mean correction in the homogeneous
area row, and the graph \eqref{JP:eq:normal-sine-graph-close} therefore
produce
$Q_\alpha\psi\in\Pi_{\rm ng}\mathcal I_\alpha Z_\alpha$ with
\begin{equation}\label{JP:eq:normal-QI}
 \|\psi-Q_\alpha\psi\|_{H^{1/2}}
 \le CS\|\psi\|_{H^{3/2}}.
\end{equation}
No overlap or ratio constant occurs.

Write the normalized disk multiplier as
\[
 \widehat q_D(f)
 =q_+(f)-k(f),\qquad
 k(f)=\sum_{|n|\ge2}k_{|n|}|\widehat f_n|^2,
 \qquad 0<k_n<\frac4n.
\]
Let $f_\alpha$ be the $q_+$-minimizing residual of
$\Pi_{\rm ng}G_\alpha^{\rm e}+W_\alpha$.  If
$g\in H_{\rm ng}^{1/2}$ and $(|D|+1)\psi=g$, Galerkin orthogonality and
\eqref{JP:eq:normal-QI} give
\[
 |\langle f_\alpha,g\rangle|
 =|q_+(f_\alpha,\psi-Q_\alpha\psi)|
 \le CS\|f_\alpha\|_{H^{1/2}}\|g\|_{H^{1/2}}.
\]
Taking the dual supremum proves the negative-norm gain
\[
 \|f_\alpha\|_{H^{-1/2}}
 \le CS\|f_\alpha\|_{H^{1/2}}.
\]
Since $q_+(f_\alpha)\le q_+(G_\alpha^{\rm e})\le CS^3$, the exact
Hilbert Schur identity
\begin{align*}
 R_B(\alpha)&:=P_{+,\sin}^{\rm e}(\alpha)-P_D^{\rm e}(\alpha)\\
 &=k(f_\alpha,f_\alpha)
  +\langle g_\alpha,
     (\widehat q_D|_{W_\alpha})^{-1}g_\alpha\rangle,
 &g_\alpha(v)&=k(f_\alpha,v),
\end{align*}
where
$P_D^{\rm e}=\inf_{v\in W_\alpha}\widehat q_D
(\Pi_{\rm ng}G_\alpha^{\rm e}+v)$, gives
\begin{equation}\label{JP:eq:normal-bessel-absolute}
 0\le R_B(\alpha)\le CS^5.
\end{equation}

For the source replacement, the exact decomposition on every side-normal
cell is
\[
 h_\alpha\sec x-1=(h_\alpha-1)+G_\alpha^{\rm e}+E_\alpha,
 \qquad E_\alpha=h_\alpha(\sec x-1)-\frac{x^2}{2}.
\]
Here $x$ is measured from the side normal, so $G_\alpha^{\rm e}=x^2/2$
on both incident half-cells.  Taylor's formula through one normalized
derivative gives
\[
 |E_\alpha|+w|E_\alpha'|\le CS^2w^2
\]
on a side cell of total angular width $w$.  Splitting the Gagliardo
integral into same, adjacent, and dyadically separated cells yields the
ratio-free estimate
\begin{equation}\label{JP:eq:normal-secant-remainder}
 \|E_\alpha\|_{H^{1/2}}
 \le CS^2\|G_\alpha^{\rm e}\|_{H^{1/2}}.
\end{equation}
The constant term $h_\alpha-1$ disappears under $\Pi_{\rm ng}$.
Coercivity of $\widehat q_D$ and the Hilbert distance inequality therefore
give
\begin{equation}\label{JP:eq:normal-secant-absolute}
 |R_{\rm sec}(\alpha)|\le CS^5,
 \qquad R_{\rm sec}=\widehat P_D^{\rm ns}-P_D^{\rm e}.
\end{equation}

We finally turn both absolute corrections into relative ones.  In a
quasiuniform material band, pull all intervals to a fixed real fan and put
$h=\|\delta\|_{\ell^\infty}/a$.  The first two material derivatives of
$q_+$ are bounded by $Ch^r$ in $H^{1/2}$, while those of $k$ are bounded by
$Ch^r$ in $H^{-1/2}$; the latter follows from
$|\Delta^rk_n|\le C_r(1+n)^{-1-r}$ and discrete summation by parts.
Differentiating the Galerkin equation and the negative-norm estimate gives,
for $r=0,1,2$,
\[
 \|\partial_t^rf_t\|_{H^{1/2}}\le Ch^ra^{3/2},
 \qquad
 \|\partial_t^rf_t\|_{H^{-1/2}}\le Ch^ra^{5/2}.
\]
Twice differentiating the displayed Schur identity gives
\[
 |D^2R_B(\gamma)[\delta,\delta]|
 \le Ca\sum_i\gamma_i^2\delta_i^2.
\]
For $R_{\rm sec}$, affine normalization of the two source profiles gives
\[
 \|\partial_t^rG_{\gamma(t)}^{\rm e}\|_{H^{1/2}}
 \le Ch^ra^{3/2},\qquad
 \|\partial_t^rE_{\gamma(t)}\|_{H^{1/2}}
 \le Ch^ra^{7/2},\qquad r=0,1,2.
\]
Differentiating the two coercive Schur equations in
\eqref{JP:eq:normal-secant-absolute} consequently gives the same Hessian
bound for $R_{\rm sec}$.  Hence their sum
$R_D=\widehat P_D^{\rm ns}-P_{+,\sin}^{\rm e}$ satisfies
\begin{equation}\label{JP:eq:normal-RD-hessian}
 |D^2R_D(\gamma)[\delta,\delta]|
 \le Ca\sum_i\gamma_i^2\delta_i^2.
\end{equation}

If $\Jfour\ge S^{9/2}$, subtract the endpoint bounds
\eqref{JP:eq:normal-bessel-absolute} and
\eqref{JP:eq:normal-secant-absolute}; this gives
$CS^{1/2}\Jfour$.  Otherwise the exact fourth-defect identity used in the
preceding lemma makes the radial segment quasiuniform.  Cyclic symmetry
removes the first derivative of $R_D$ at the regular fan, and Taylor's
formula with \eqref{JP:eq:normal-RD-hessian} gives $CS\Jfour$.  This proves
\eqref{JP:eq:normal-bessel-secant-relative} with
$\omega_D(S)=C(S^{1/2}+S)$.
\end{proof}

\begin{theorem}[Closed fan-normal disk endpoint and principal forcing]
\label{JP:thm:normal-disk-endpoint}
There are universal $S_D,\kappa,C_p>0$ such that every positive fan with
$S\le S_D$ satisfies
\begin{align}
 P_D^{\rm ns}(\alpha)-P_D^{\rm ns}(\alpha_*)
 &\ge2\kappa\Jfour(\alpha),
 \label{JP:eq:D0ns}\\
 \|p_\alpha\|_{(\Gjp_\alpha)^*}
 &\le C_p\Jfour(\alpha)^{1/2}.
 \label{JP:eq:JE-P}
\end{align}
Both constants are independent of all minimum angles and adjacent ratios.
\end{theorem}

\begin{proof}
Combine \eqref{JP:eq:normal-homogeneous-gap},
\eqref{JP:eq:normal-sine-relative}, and
\eqref{JP:eq:normal-bessel-secant-relative}, then decrease $S_D$ so that the
two moduli consume at most half of the homogeneous margin.  Multiplication
by $2\pi d^2$ gives \eqref{JP:eq:D0ns} for a universal $\kappa>0$.

The closed exact source envelope from F129 gives
\[
 q_D(B_\alpha)-q_D(B_{\alpha_*})\le C\Jfour(\alpha).
\]
Hilbert completion and $p_{\alpha_*}=0$ give
\[
 \frac14\|p_\alpha\|_{(\Gjp_\alpha)^*}^2
 =\{q_D(B_\alpha)-q_D(B_{\alpha_*})\}
  -\{P_D^{\rm ns}(\alpha)-P_D^{\rm ns}(\alpha_*)\}.
\]
Drop the favorable second brace and take square roots to obtain
\eqref{JP:eq:JE-P}.
\end{proof}

\begin{theorem}[Closed exact fan-normal value matching]
\label{JP:thm:normal-value-matching}
There is a modulus $\varepsilon_V(S)\to0$, independent of all minimum
angles and adjacent ratios, such that every sufficiently fine positive fan
satisfies
\begin{equation}\label{JP:eq:JE-V}
 \boxed{
 \left|\Delta_\alpha-
 \{q_D(B_\alpha)-q_D(B_{\alpha_*})\}\right|
 \le\varepsilon_V(S)\Jfour(\alpha).}
\end{equation}
\end{theorem}

\begin{proof}
Let $W_\alpha=B_\alpha e_r$ be the radial physical boundary field of the
area-normalized tangential polygon; its normal scalar trace at the disk is
the source $B_\alpha$ in \eqref{JP:eq:disk-source-B}.  The exact
disk-normalized scalar expansion used in JP--BU may be retained
at its natural orders rather than weakened to a one-sided bound.  It is
\begin{equation}\label{JP:eq:value-matching-split}
 \Delta_\alpha
 =\{q_D(B_\alpha)-q_D(B_{\alpha_*})\}
  +\{\Delta_3(\alpha)-\Delta_3(\alpha_*)\}
  +\{F_4(\alpha)-F_4(\alpha_*)\},
\end{equation}
where $\Delta_3$ is the complete scalar cubic row, including its factor
$1/6$, and $F_4$ is the exact fourth-and-higher scalar remainder.

The closed CT cubic identity splits $\Delta_3$ into its principal null form
and the complete normal/geometric order-one remainder.  CT's relative
closure theorem controls the latter by $o(\Jfour)$.  The QB
secant-to-quadratic transfer and the closed DF theorem control the principal
null form by the same scale.  Consequently there is a modulus
$\varepsilon_3(S)\to0$ with
\begin{equation}\label{JP:eq:value-cubic-relative}
 |\Delta_3(\alpha)-\Delta_3(\alpha_*)|
 \le\varepsilon_3(S)\Jfour(\alpha).
\end{equation}
This is the pure scalar row, so no support-metric comparison is involved.

It remains to sharpen the tail estimate already used in JP--BU.  On the
real chart and its material complex disk, the tame scalar remainder gives
\begin{equation}\label{JP:eq:value-tail-absolute}
 |F_4(\zeta)|
 \le C\|W_\zeta\|_{H^{1/2}}^2
       \|W_\zeta\|_{W^{1,\infty}}^2
 \le CS^5.
\end{equation}
Put $d_i=\alpha_i-a$ and $\gamma_i(t)=a+td_i$.  If
$\Jfour\ge S^{9/2}$, applying \eqref{JP:eq:value-tail-absolute} at the two
endpoints gives
\[
 |F_4(\alpha)-F_4(\alpha_*)|
 \le CS^{1/2}\Jfour(\alpha).
\]
If $\Jfour<S^{9/2}$, the exact fourth-defect identity implies
$a\asymp S$ and keeps the whole radial segment quasiuniform.  The material
complex radius in direction $d$ is at least
$ca/\|d\|_{\ell^\infty}$.  Cauchy's estimate and
\eqref{JP:eq:value-tail-absolute} give, with
$E(t)=\sum_i\gamma_i(t)^2d_i^2$,
\[
 |F_4''(t)|
 \le CS^5\frac{\|d\|_{\ell^\infty}^2}{a^2}
 \le CS E(t).
\]
Cyclic symmetry makes $F_4'(0)$ constant in the fan coordinates and hence
$F_4'(0)[d]=0$.  Taylor integration and
$\Jfour=12\int_0^1(1-t)E(t)\,dt$ give the near estimate
$CS\Jfour$.  Together with \eqref{JP:eq:value-cubic-relative} and
\eqref{JP:eq:value-matching-split}, this proves \eqref{JP:eq:JE-V} with
$\varepsilon_V(S)=\varepsilon_3(S)+C(S^{1/2}+S)$.
\end{proof}

\begin{proposition}[Exact forcing-residual covector]
\label{JP:prop:forcing-residual-split}
On the polar side cell $K_i$, put
\[
 f_\alpha(\theta)=h_\alpha\sec(\theta-\theta_i)-1
\]
and let
$Z_\alpha(z)=v_\alpha(z)e_r+w_\alpha(z)e_\theta$ be the continuous
physical-affine support displacement.  In the notation of the complete CT
tensor, let $\kappa_{\mathbb D}\mathcal C+\mathcal R_{\rm full}$ be the
normal third-derivative split and set
\[
 \Psi_{4,\alpha}[z]
 :=D\mathcal R_4(f_\alpha e_r)[Z_\alpha(z)].
\]
Then the exact residual $e_\alpha=\ell_\alpha-p_\alpha$ is
\begin{equation}\label{JP:eq:exact-forcing-residual}
\boxed{
\begin{aligned}
 e_\alpha[z]={}&
 2\langle f_\alpha,
       v_\alpha(z)-\mathcal I_\alpha z\rangle_D\\
 &+\frac{\kappa_{\mathbb D}}2
   \mathcal C(f_\alpha,f_\alpha,v_\alpha(z))
  +\frac12\mathcal R_{\rm full}
   (f_\alpha,f_\alpha,v_\alpha(z))\\
 &-2\langle f_\alpha,
       w_\alpha(z)f_\alpha'\rangle_D
  +\Psi_{4,\alpha}[z].
\end{aligned}}
\end{equation}
Every term is evaluated after the exact area/translation projection.  The
formula is an identity of covectors on $\mathsf X_\alpha$; in particular,
it contains no norm comparison and no asymptotic remainder notation.
\end{proposition}

\begin{proof}
The boundary of the tangential polygon is the disk boundary displaced by
$W_\alpha=f_\alpha e_r$.  Because vertices and side parametrizations are
affine in the support numbers, differentiation in the support direction is
the single physical slot $Z_\alpha(z)$.  The exact scalar Taylor identity
therefore gives
\[
 \ell_\alpha[z]
 =D^2\mathcal L(0)[W_\alpha,Z_\alpha(z)]
 +\frac12D^3\mathcal L(0)
       [W_\alpha,W_\alpha,Z_\alpha(z)]
 +\Psi_{4,\alpha}[z].
\]
The disk Hessian is
$D^2\mathcal L(0)[W_\alpha,Z_\alpha(z)]
=2\langle f_\alpha,v_\alpha(z)\rangle_D$.
The complete physical-coordinate CT identity gives
\[
\begin{aligned}
 \frac12D^3\mathcal L(0)
 [W_\alpha,W_\alpha,Z_\alpha(z)]
 ={}&\frac{\kappa_{\mathbb D}}2
   \mathcal C(f_\alpha,f_\alpha,v_\alpha(z))\\
 &+\frac12\mathcal R_{\rm full}
   (f_\alpha,f_\alpha,v_\alpha(z))
 -2\langle f_\alpha,w_\alpha(z)f_\alpha'\rangle_D.
\end{aligned}
\]
The last coefficient is the curvature multiplicity after differentiating
the cubic Taylor coefficient: the raw third derivative contributes
$-4\langle f_\alpha,w_\alpha(z)f_\alpha'\rangle_D$ and is multiplied by
$1/2$.  Finally
$p_\alpha[z]=2\langle f_\alpha,\mathcal I_\alpha z\rangle_D$ because the
disk form annihilates the constant and first harmonics.  Subtraction proves
\eqref{JP:eq:exact-forcing-residual}.
\end{proof}

\begin{remark}[Retraction of the period-three material forcing transfer]
\label{JP:rem:period-three-material-forcing}
The former proposition which identified the three-cell material constant
$R_3$ with the genuine polygonal covector
$e_\alpha=\ell_\alpha-p_\alpha$ is withdrawn.  The revision-22
Schwarz--Christoffel audit finds a same-order corner finite part already in
the support diagonal, so the material-to-polygon transfer used below does
not justify the mixed forcing coefficient either.  The displayed material
series remains a useful diagnostic, but it proves neither truth nor
falsity of the former JE--FM estimate
\begin{equation}\label{JP:eq:JE-FM}
 \|e_\alpha\|_{(\mathfrak G_\alpha^{\rm ns})^*}
 \le\varepsilon(S)\Jfour(\alpha)^{1/2},
 \qquad \varepsilon(S)\longrightarrow0.
\end{equation}
JE--FM is absent from the active
DAG, and no active implication uses this archived calculation.
\end{remark}

\iffalse
For a three-periodic material function define
\[
 \widehat h(n)=\frac13\int_0^3h(y)e^{-2\pi iny/3}\,dy,
 \qquad
 \langle f,g\rangle_3
 =\sum_{n\ne0}\left|\frac{2\pi n}{3}\right|
       \widehat f(-n)\widehat g(n).
\]
On a normal cell of length $l_i$, let $G_l$ be the edge quadratic
source, let $S_Z$ equal $Z_i$ and $Z_{i+1}$ on the two half-cells, and let
$I_l^0Z$ be the linear interpolant.  If $\varphi_i$ is the cell midpoint,
direct integration gives, for $\omega=2\pi n/3\ne0$,
\begin{align}
 \widehat G_l(\omega)
 &=\frac1{3\omega^2}\sum_i l_i e^{-i\omega\varphi_i},\label{JP:eq:G-Fourier}\\
 \widehat{(S_Z-I_l^0Z)}(\omega)
 &=\frac1{3i\omega}\sum_i(Z_{i+1}-Z_i)e^{-i\omega\varphi_i}
 \left(1-\operatorname{sinc}\frac{\omega l_i}{2}\right).
 \label{JP:eq:step-sine-Fourier}
\end{align}
We now give the alias calculation, including the realification.  Put
\[
 \kappa=\frac{2\pi}{3},\qquad
 D_i^+=e^{i\kappa i},\qquad Z_i^-=e^{-i\kappa i},
\]
\[
 \Delta=e^{-i\kappa}-1,\qquad
 A=\frac1{e^{i\kappa}-1}+\frac12=-\frac{i}{2\sqrt3}.
\]
For the complexified material family,
\[
 l_i(t)=1+tD_i^+,\qquad
 \varphi_i(t)=i+\frac12+tAD_i^+.
\]
Let $H_t=S_{Z^-}-I_{l(t)}^0Z^-$ and
$\omega_m=2\pi(m-\frac13)$.  Direct substitution in
\eqref{JP:eq:G-Fourier}--\eqref{JP:eq:step-sine-Fourier} gives
\begin{align}
 \widehat H_0(\omega_m)
 &=\frac{\Delta e^{-i\omega_m/2}}{i\omega_m}
 \left(1-\operatorname{sinc}\frac{\omega_m}{2}\right),
 \label{JP:eq:R3-H-first}\\
 \left.\partial_t\right|_0\widehat G_t(-\omega_m)
 &=\frac{e^{i\omega_m/2}}{\omega_m^2}
 \left(1+\frac{\omega_m}{2\sqrt3}\right).
 \label{JP:eq:R3-G-first}
\end{align}
The second alias family comes from the regular source and the material
derivative of $H_t$.  For $m\ne0$,
\begin{align}
 \widehat G_0(2\pi m)&=\frac{(-1)^m}{(2\pi m)^2},
 \label{JP:eq:R3-G-second}\\
 \left.\partial_t\right|_0\widehat H_t(2\pi m)
 &=\frac{\Delta(-1)^m}{i}
 \left(-\frac1{2\sqrt3}-\frac{(-1)^m}{2\pi m}\right).
 \label{JP:eq:R3-H-second}
\end{align}
Consequently the complex bilinear derivative is
\begin{equation}\label{JP:eq:R3-alias-complex}
 \left.\partial_t\right|_0\langle G_t,H_t\rangle_3
 =\frac{\Delta}{i}R_3,
\end{equation}
where the two sums below are taken with the same symmetric cutoff:
\begin{align}
 R_3={}&
 \sum_{m\in\mathbb Z}
 \frac{\operatorname {sgn}\omega_m}{\omega_m^2}
 \left(1+\frac{\omega_m}{2\sqrt3}\right)
 \left(1-\operatorname{sinc}\frac{\omega_m}{2}\right)
 \notag\\
 &+\sum_{m\ne0}\frac1{2\pi|m|}
 \left(-\frac1{2\sqrt3}-\frac{(-1)^m}{2\pi m}\right).
 \label{JP:eq:R3-alias-sum}
\end{align}
This joint cutoff is essential: the logarithmic terms cancel between the
two alias families.

For $x_m=m-\frac13$,
\[
 1-\operatorname{sinc}\frac{\omega_m}{2}
 =1+\frac{(-1)^m\sqrt3}{2\pi x_m}.
\]
Expanding \eqref{JP:eq:R3-alias-sum} gives
\begin{equation}\label{JP:eq:R3-four-sums}
 R_3=\frac{S_1}{4\pi^2}
 +\frac{S_0}{4\pi\sqrt3}
 +\frac{\sqrt3S_3}{8\pi^3}
 +\frac{S_2}{8\pi^2},
\end{equation}
where, with the same symmetric convention,
\[
\begin{aligned}
 S_1&=\sum_{m\in\mathbb Z}
 \frac{\operatorname {sgn}x_m}{x_m^2},&
 S_0&=\sum_{m\in\mathbb Z}\frac1{|x_m|}
      -\sum_{m\ne0}\frac1{|m|},\\
 S_3&=\sum_{m\in\mathbb Z}
 \frac{(-1)^m\operatorname {sgn}x_m}{x_m^3},&
 S_2&=\sum_{m\in\mathbb Z}
 \frac{(-1)^m\operatorname {sgn}x_m}{x_m^2}.
\end{aligned}
\]
The remaining odd summand satisfies
$\sum_{m\ne0}(-1)^m/(|m|m)=0$.  Pairing the positive and negative
indices and then grouping by residue classes modulo $3$ and $6$ yields
\begin{equation}\label{JP:eq:R3-residue-sums}
 S_1=-9L(2,\chi_3),\qquad S_0=3\log3,\qquad
 S_3=\frac{39}{2}\zeta(3),\qquad
 S_2=-\frac{27}{2}L(2,\chi_3).
\end{equation}
Equations \eqref{JP:eq:R3-four-sums}--\eqref{JP:eq:R3-residue-sums} give
\eqref{JP:eq:R3-closed}.

The real family in the proposition is obtained from
\[
 D=\operatorname {Re}D^+,\qquad
 Z=\operatorname {Re}(e^{-i\kappa}Z^-).
\]
Real bilinearity and \eqref{JP:eq:R3-alias-complex} therefore give the
separate realification identity
\begin{equation}\label{JP:eq:R3-alias-real}
 \left.\partial_t\right|_0
 \langle G_{\alpha(t)},S_Z-I_{\alpha(t)}^0Z\rangle_3
 =\frac12\operatorname {Re}
 \left(e^{-i\kappa}\frac{\Delta}{i}\right)R_3
 =\frac{\sqrt3}{2}R_3.
\end{equation}

The sign in \eqref{JP:eq:R3-closed} is rigorous: write
\[
 L(2,\chi_3)=\sum_{k\ge0}
 \frac{6k+3}{(3k+1)^2(3k+2)^2}.
\]
After the first $1000$ terms its tail lies between zero and
$1/(18\cdot999^2)$.  Integral tails after $1000$ terms for $\zeta(3)$,
the positive series
\[
\log3=2\sum_{j\ge0}\frac{2^{-(2j+1)}}{2j+1},
\]
the rational enclosures
\[
 \frac{333}{106}<\pi<\frac{355}{113},\qquad
 \frac{1732050807}{10^9}<\sqrt3<\frac{1732050808}{10^9}
\]
give the displayed rational interval for $R_3$.

On each side cell, in the material coordinate
$x\in[-1/2,1/2]$, one has, also after one $t$ derivative,
\begin{equation}\label{JP:eq:R3-cell-expansions}
\begin{gathered}
 f_\alpha=a^2\left(\frac{x^2}{2}-\frac1{24}\right)+O(a^4),\\
 v_\alpha(Z)-w_\alpha(Z)f_\alpha'=S_Z+O(a^2),
 \qquad \mathcal I_\alpha Z=I^0Z+O(a^2).
\end{gathered}
\end{equation}
The corresponding material norms satisfy
\begin{equation}\label{JP:eq:R3-material-scales}
 \|f_\alpha\|_\infty=O(a^2),\quad
 \|f_\alpha'\|_\infty=O(a),\quad
 \|f_\alpha\|_{H^{1/2}}=O(a^{3/2}),\quad
 \|v_\alpha(Z)\|_{H^{1/2}}=O(a^{-1/2}).
\end{equation}
Thus the unique $at$ term in
\eqref{JP:eq:exact-forcing-residual} is
$2q_D(f_\alpha,S_Z-I^0Z)$.  Indeed, the high-frequency symbol
\[
 \mathcal D_\circ+1=|D|+1+O(|D|^{-1})
\]
and the $2\pi/(3a)$ material periods show that the $|D|$ row is
$4\pi d_\circ^2at$ times \eqref{JP:eq:R3-alias-real}; the $+1$ row is
$O(a^2|t|)$, while the order-minus-one row and the errors in
\eqref{JP:eq:R3-cell-expansions} are $O(a^3|t|)$.

For completeness, the one-Lipschitz null-form estimate gives
\[
 |\mathcal C(f,f,v)|
 \le C\|f'\|_\infty
       \|f\|_{H^{1/2}}\|v\|_{H^{1/2}}=O(a^2).
\]
At $t=0$ the period-one source is Fourier-orthogonal to the period-three
mode, so along the family
\begin{equation}\label{JP:eq:R3-lower-orders}
 \mathcal C(f_t,f_t,v_t)=O(a^2|t|),\qquad
 \mathcal R_{\rm full}(f_t,f_t,v_t)=O(a^3|t|).
\end{equation}
We record the special-family fourth-tail estimate rather than invoking
the retracted general support-covector row CT--M4.  Write the analytic
disk expansion as
\[
 \Psi_{4,\alpha(t)}[Z]
 =\sum_{k\ge4}\frac1{(k-1)!}
 D^k\mathcal L(0)
 [f_{\alpha(t)}^{\,k-1},Z_{\alpha(t)}(Z)].
\]
After rescaling each of the three material cells to unit length,
$f=a^2F_{a,t}$ and $f'=a\,\partial_xF_{a,t}$, with the profiles
holomorphic and uniformly bounded in $t$ on a fixed complex
neighbourhood.  On the
physical circle the relevant full-vector bounds are
\begin{equation}\label{JP:eq:R3-full-vector-scales}
\begin{gathered}
 \|f_{\alpha(t)}e_r\|_{H^{1/2}}\le Ca^{3/2},\qquad
 \|f_{\alpha(t)}e_r\|_{W^{1,\infty}}\le Ca,\\
 \|Z_{\alpha(t)}(Z)\|_{H^{1/2}}\le Ca^{-3/2}.
\end{gathered}
\end{equation}
To see the last scale, the normal component has period $3a$ and size one,
while the
tangential endpoint coefficients are $O(a^{-1})$; the piecewise-affine
full vector therefore has $H^{1/2}$ size $O(a^{-3/2})$.  We do not place
this support slot in $W^{1,\infty}$, where its derivative is
$O(a^{-2})$.

The fixed-chart analytic Hessian estimate is bilinear in two
$H^{1/2}$ slots and analytic in the remaining
$W^{1,\infty}$ slots.  Applying its Cauchy bound with one copy of $f$
and the support direction in the two Hilbert slots gives, uniformly for
$|t|\le1/2$,
\begin{equation}\label{JP:eq:R3-special-tail-bound}
 \frac1{(k-1)!}
 \left|D^k\mathcal L(0)
 [f_{\alpha(t)}^{\,k-1},Z_{\alpha(t)}(Z)]\right|
 \le C_0C_1^k
 \|f e_r\|_{H^{1/2}}\|Z_\alpha(Z)\|_{H^{1/2}}
 \|f e_r\|_{W^{1,\infty}}^{k-2}
 \le C_0C_1^k a^{k-2}.
\end{equation}
For $C_1a<1/2$, summing \eqref{JP:eq:R3-special-tail-bound} gives
$|\Psi_4(t)|\le Ca^2$.

The derivative in $t$ is taken only after passing to the fixed
three-cell material representation.  There the side-to-material maps,
Jacobians, and tensor coefficients are rational holomorphic functions of
$t$ on a disk $|t|<r$ with $r>0$ independent of $a$.  Thus the
\emph{scalar} function $t\mapsto\Psi_4(t)$ is holomorphic on that disk,
and the preceding $Ca^2$ bound holds uniformly there.  One-variable
Cauchy gives $|\Psi_4'(t)|\le C_ra^2$ for $|t|<r/2$.  This argument does
not differentiate a moving-knot physical field in $H^{1/2}$.
The value at $t=0$ vanishes by cyclic Fourier orthogonality, hence
\[
 \Psi_{4,\alpha(t)}[Z]=O(a^2|t|).
\]
The nonlinear material geometry contributes $O(at^2)$.  Combining these
estimates with \eqref{JP:eq:R3-alias-real} first proves the asserted
expansion for the unprojected $Z$.

It remains to check the exact quotient projection.  From
\eqref{JP:eq:scale-section} and the tangential side lengths,
\begin{align}
 s_{\alpha(t)}(Z)
 &=\frac{h_{\alpha(t)}}{2\pi}
 \sum_i\tan\frac{\alpha_i(t)}2\,(Z_i+Z_{i+1})\notag\\
 &=-\frac{h_{\alpha(t)}}4t+O(a^2|t|+t^2).
 \label{JP:eq:R3-scale-projection}
\end{align}
Thus, for a translation vector $\tau_{\alpha(t)}$,
\begin{equation}\label{JP:eq:R3-projected-Z}
 Z_{\alpha(t)}
 =Z+\frac t4\mathbf1+\tau_{\alpha(t)}
   +O(a^2|t|+t^2)\mathbf1.
\end{equation}
Both covectors annihilate translations, and scale invariance gives
$\ell_\alpha[\mathbf1]=0$; in contrast, $p_\alpha[\mathbf1]$ need not
vanish.  On the present quasiuniform family the material-cell expansions
\[
 B_{\alpha(t)}=a^2F_{a,t}+O(a^4),\qquad
 K_{\alpha(t)}\mathbf1=a^2C_{a,t}+O(a^4)
\]
have uniformly bounded three-periodic profiles.  Since $q_D$ has
material order one,
\begin{equation}\label{JP:eq:R3-scale-p}
 |p_{\alpha(t)}[\mathbf1]|
 =2|\langle B_{\alpha(t)},K_{\alpha(t)}\mathbf1\rangle_D|
 \le Ca^3.
\end{equation}
Equations \eqref{JP:eq:R3-projected-Z}--\eqref{JP:eq:R3-scale-p} change
$e_\alpha[Z]$ by only $O(a^3|t|+a^3t^2)$, which is contained in
$O(a^2|t|+at^2)$.  This proves \eqref{JP:eq:period-three-e} for the exact
projected direction.

Finally
\begin{align*}
 \Jfour(\alpha(t))&=6\pi a^3t^2+O(a^3|t|^3),\\
 \mathfrak G_{\alpha(t)}^{\rm ns}(Z_{\alpha(t)})
 &=\frac{2\pi d_\circ^2}{a}Q_3+O(1),
 \qquad Q_3=\langle I^0Z,I^0Z\rangle_3>0.
\end{align*}
Hence the quotient of the left side of \eqref{JP:eq:period-three-e} by
$\Jfour^{1/2}(\mathfrak G^{\rm ns})^{1/2}$ tends in absolute value to
$|d_\circ|R_3/\sqrt{Q_3}>0$.  Taking $t=a$ keeps all angles positive and
gives $\Jfour\asymp a^5<S^{9/2}$, so the counterexample lies in the near
branch and has adjacent ratios tending to one.
\end{proof}
\fi

The closed disk Schur, forcing, and value rows are retained as JP--DE.
The unproved separated flux interface is absent from the active DAG; this
does not remove the fan-normal metric.  The replacement is the fixed-loss
joint scalar comparison in
\texttt{JC\_joint\_comparison.tex}.

\begin{remark}[Retraction of the period-three joint-loss transfer]
\label{JP:rem:period-three-material-joint-loss}
The former period-three joint-loss proposition is also withdrawn.  Its
coefficient $E_3$ is the material cell value, while the true polygonal
support finite part is $P(1/3)/2=9\log3/(4\pi)$.  Moreover the material
constant $R_3$ has not been transferred through the mixed polygonal corner
layer.  Consequently the limit formerly labelled
\textup{``joint-period3-loss''} is not a theorem about the genuine
polygonal residual.  The active JC--PL formulation already uses a fixed
absorbable loss and does not depend on this archived diagnostic.
\end{remark}

\iffalse
The exact regular-fan leading terms furnished by the displayed material
profiles are
\begin{align*}
 D_h^2\Phi_{\alpha_*}[Z,Z]
 &=\frac{4\pi d_\circ^2}{a}E_3+O(1),\\
 \mathfrak G_{\alpha_*}^{\rm ns}(Z)
 &=\frac{2\pi d_\circ^2}{a}Q_3+O(1).
\end{align*}
Here is an explicit certificate for \eqref{JP:eq:E3Q3-bounds}.  Put
\[
 \omega_n=\frac{2\pi n}{3},\qquad
 J_k(\omega)=\int_{-1/2}^{1/2}x^ke^{-i\omega x}\,dx.
\]
Then
\begin{align}
 J_0(\omega)&=\frac{2\sin(\omega/2)}{\omega},\label{JP:eq:J0-recurrence}\\
 J_k(\omega)&=
 -\frac{2^{-k}e^{-i\omega/2}-(-2)^{-k}e^{i\omega/2}}{i\omega}
 +\frac{k}{i\omega}J_{k-1}(\omega).\label{JP:eq:Jk-recurrence}
\end{align}
For a cell polynomial $P_r(x)=\sum_{j=0}^2p_{rj}x^j$,
\begin{equation}\label{JP:eq:poly-Fourier}
 \widehat P(n)=\frac13\sum_{r=0}^2e^{-i\omega_nr}
 \sum_{j=0}^2p_{rj}J_j(\omega_n).
\end{equation}
The three cell polynomials are
\[
\begin{array}{c|ccc}
r&0&1&2\\ \hline
V_r&-\frac12-\frac34x+\frac32x^2&
 -\frac12+\frac34x+\frac32x^2&1-3x^2\\
U_r&-\frac34x+\frac32x^2&
 \frac34x+\frac32x^2&-3x^2\\
K_r&-\frac9{16}+\frac92x-\frac{27}{4}x^2&
 -\frac9{16}-\frac92x-\frac{27}{4}x^2&-27x^2.
\end{array}
\]
For the linear interpolant, put
\[
 L_0(\omega)=\frac{1-e^{-i\omega}}{i\omega},\qquad
 L_1(\omega)=-\frac{e^{-i\omega}}{i\omega}
              +\frac{L_0(\omega)}{i\omega};
\]
then
\begin{equation}\label{JP:eq:I-Fourier}
 \widehat I(n)=\frac13\sum_{r=0}^2e^{-i\omega_nr}
 \{Z_rL_0(\omega_n)+(Z_{r+1}-Z_r)L_1(\omega_n)\}.
\end{equation}
Define
\[
\begin{aligned}
 A_3&=\sum_{n\ne0}|\omega_n||\widehat V(n)|^2,&
 B_3&=\sum_{n\ne0}|\omega_n|
       \overline{\widehat F(n)}\widehat K(n),\\
 C_3&=\sum_{n\ne0}|\omega_n|
       \overline{\widehat V(n)}\widehat U(n),&
 Q_3&=\sum_{n\ne0}|\omega_n||\widehat I(n)|^2.
\end{aligned}
\]
Thus $E_3=A_3+B_3-2C_3$.  Symmetric summation over
$0<|n|\le4096$, using
\eqref{JP:eq:J0-recurrence}--\eqref{JP:eq:I-Fourier}, gives the outward
intervals
\[
\begin{aligned}
0.76393070&<A_K<0.76393075,&
-0.26168518&<B_K<-0.26168512,\\
-0.13865959&<C_K<-0.13865952,&
0.56698459&<Q_K<0.56698465.
\end{aligned}
\]
Only $0,\pm1/2,\pm\sqrt3/2,\pm1$ occur in the trigonometric factors, so
these are finite rational interval calculations using rational outward
bounds for $\pi$ and $\sqrt3$.  Finally
\[
\operatorname {TV}(F')=6,\quad \operatorname {TV}(V')=15,\quad
\operatorname {TV}(I')=6,\quad
\operatorname {TV}(K)\le\frac{135}{4},\quad
\operatorname {TV}(U)\le\frac{15}{2}.
\]
Twice integrating the continuous profiles and once integrating the BV
profiles gives
\[
\begin{aligned}
0\le A_3-A_K&<1.63\cdot10^{-7},&
0\le Q_3-Q_K&<2.60\cdot10^{-8},\\
|B_3-B_K|&<0.002505,&
|C_3-C_K|&<0.001392.
\end{aligned}
\]
These intervals imply
$0.56698456<Q_3<0.56698468$ and
$0.2072<E_3-Q_3<0.2179$, hence
\eqref{JP:eq:E3Q3-bounds}.

We finish by checking the projection, active interval, and every Taylor
remainder on the stated path $t=a$.  Equations
\eqref{JP:eq:R3-projected-Z}--\eqref{JP:eq:R3-scale-p}, the analytic
three-cell material pullback, and the two leading terms at the start of
the proof give, uniformly for $|s|\le c_0a^2$,
\begin{align}
 \mathscr R_{\alpha(a)}(sZ_{\alpha(a)})
 ={}&
 \{2\pi\sqrt3\,d_\circ^2R_3a^2+O(a^3)\}s\notag\\
 &+\left\{\frac{2\pi d_\circ^2}{a}(E_3-Q_3)+O(1)\right\}s^2
 +O(a^{-2}|s|^3).
 \label{JP:eq:joint-period3-Taylor}
\end{align}
Here the $O(1)s^2$ term includes the material change
$H_{\alpha(a)}-H_{\alpha_*}$ and the scale projection; the latter uses
$K_\alpha\mathbf1=O(a^2)$ in material amplitude.  The last term follows
by contracting the third support derivative on the fixed three material
cells: its circular multipliers have total order at most two, hence its
norm on three-periodic support data is $O(a^{-2})$.

Side lengths are affine in the support numbers.  Their base size is
comparable to $a$, while a support increment $sZ_{\alpha(a)}$ changes
them by $O(|s|/a)$.  Choosing $c_0$ sufficiently small therefore keeps
the entire interval $|s|\le c_0a^2$ active.  The minimizer of the
quadratic part of \eqref{JP:eq:joint-period3-Taylor} is
\begin{equation}\label{JP:eq:joint-period3-minimizer}
 s_*=-\frac{\sqrt3R_3}{2(E_3-Q_3)}a^3+O(a^4),
\end{equation}
so $|s_*|\ll a^2$.  At this point the linear-coefficient error,
the $O(1)s^2$ term, the projection error, and the support cubic remainder
are respectively $O(a^6)$, $O(a^6)$, $O(a^6)$, and $O(a^7)$.
Since
\[
 \Jfour(\alpha(a))=6\pi a^5+O(a^6),
\]
substitution of \eqref{JP:eq:joint-period3-minimizer} into
\eqref{JP:eq:joint-period3-Taylor} gives the upper bound in
\eqref{JP:eq:joint-period3-loss}.  For the matching lower bound, throughout
$|s|\le c_0a^2$ the cubic term divided by the positive leading quadratic
term is $O(|s|/a)\le O(c_0a)$.  Hence outside
$|s|\le Ma^3$, with $M$ fixed sufficiently large, the positive quadratic
term dominates both the linear term and every remainder.  On the
remaining interval put $s=a^3\sigma$ and divide
\eqref{JP:eq:joint-period3-Taylor} by $a^5$.  The result converges uniformly
on bounded $\sigma$ to
\[
 2\pi\sqrt3\,d_\circ^2R_3\sigma
 +2\pi d_\circ^2(E_3-Q_3)\sigma^2.
\]
Its minimum, divided by
$\Jfour(\alpha(a))/a^5\to6\pi$, is exactly the right side of
\eqref{JP:eq:joint-period3-loss}.  This proves the matching lower bound and
the asserted limit of the local infimum.
\end{proof}
\fi

\begin{remark}[Why the closed TR/DM metric cannot be transferred]
TR completes its square in
$\widetilde G_\alpha(z)=q_D(\Pi_{\rm ng}\widetilde T_\alpha z)$, not in
$\Gjp_\alpha$.  These metrics have no no-ratio comparison in the needed
direction.  Indeed, take adjacent angles $\varepsilon\ll a$ and support
data $z_0=0$, $z_1=z_2=1$, return to zero across ordinary cells, and remove
the scale and translation modes.  The sine trace has only one sharp
transition, so
\[
 \Gjp_\alpha(z)\le C\{1+\log(a/\varepsilon)\}.
\]
On the adjacent physical cell, however,
 $a_1^-=(\sin\varepsilon)^{-1}-\tan(\varepsilon/2)\asymp\varepsilon^{-1}$.
On a subinterval of length comparable to $a$ this gives
$|\widetilde T_\alpha z|\ge ca/\varepsilon$.  Removing three low modes
does not remove this localized mass, and hence
\[
 \widetilde G_\alpha(z)\ge c\frac{a^3}{\varepsilon^2},
 \qquad
 \frac{\widetilde G_\alpha(z)}{\Gjp_\alpha(z)}\longrightarrow\infty.
\]
Thus the closed DM estimate in the old physical-radial metric cannot imply
the fan-normal dual estimate \eqref{JP:eq:JE-FM}.  This is a strict interface
mismatch, not a missing citation.
\end{remark}

\begin{remark}[Diagnostic check of the new Schur sign]
The reproducible script
\texttt{tests/check\_jp\_joint\_endpoint.py} samples the exact
finite-dimensional Schur complement in \eqref{JP:eq:D0ns-definition} on the
area/translation quotient.  For regular-nearby, alternating-tiny,
one-tiny-angle, and fixed-seed random fans with
$N=8,12,16,24$, the sampled ratio
\[
 \frac{P_D^{\rm ns}(\alpha)-P_D^{\rm ns}(\alpha_*)}
      {\Jfour(\alpha)}
\]
lies between $0.0189$ and $0.0333$; the sampled metric remains positive on
the quotient.  This independently checks the sign and normalization used
in Theorem~\ref{JP:thm:normal-disk-endpoint}; the script itself is not a proof
and its finite resolution does not test the mixed polygonal corner finite
part isolated in the JC node.
\end{remark}

\begin{proposition}[Archived conditional JP--FS algebra]
\label{JP:prop:FS}
Assume JP--SC with $c=2-\delta$, $0<\delta\le1/2$, and JP--JE.  After
decreasing the common fan and chart thresholds,
\begin{equation}\label{JP:eq:FS}
 \boxed{
 \Delta_\alpha
 -\frac1{2c}\|\ell_\alpha\|_{(\Gjp_\alpha)^*}^2
 \ge\kappa\Jfour(\alpha).}
\end{equation}
Equivalently, with $\mathcal Q_\alpha=c\Gjp_\alpha$, the dual term is
\(
 \frac12\sup_{z\ne0}\frac{(\sum_i r_{\alpha,i}z_i)^2}
 {\mathcal Q_\alpha(z)}
\).
\end{proposition}

\begin{proof}
Write $\|\cdot\|_* =\|\cdot\|_{(\Gjp_\alpha)^*}$ and
$e_\alpha=\ell_\alpha-p_\alpha$, and enlarge the two moduli to the common
one
\[
 \varepsilon(S):=\max\{\varepsilon_V(S),
                         \varepsilon_{\rm JE}(S)\}\longrightarrow0.
\]
Hilbert completion and
$p_{\alpha_*}=0$ give
\[
 P_D^{\rm ns}(\alpha)-P_D^{\rm ns}(\alpha_*)
 =q_D(B_\alpha)-q_D(B_{\alpha_*})
  -\frac14\|p_\alpha\|_*^2.
\]
For $c=2-\delta$ and $\delta\le1/2$,
\begin{align*}
 \frac1{2c}\|\ell_\alpha\|_*^2-\frac14\|p_\alpha\|_*^2
 &\le\left\{
 \frac{\delta C_p^2}{6}
 +\frac{2C_p\varepsilon(S)
        +\varepsilon(S)^2}{3}\right\}\Jfour(\alpha).
\end{align*}
Indeed,
$1/(2c)-1/4\le\delta/6$, $1/(2c)\le1/3$, and then use
\eqref{JP:eq:JE-P}--\eqref{JP:eq:JE-FM} and Cauchy--Schwarz.  Combining this
with \eqref{JP:eq:D0ns} and \eqref{JP:eq:JE-V} yields
\begin{align*}
 \Delta_\alpha-\frac1{2c}\|\ell_\alpha\|_*^2
 \ge\Bigl[2\kappa-\varepsilon(S)
 -\frac{\delta C_p^2}{6}
 -\frac{2C_p\varepsilon(S)
        +\varepsilon(S)^2}{3}\Bigr]\Jfour(\alpha).
\end{align*}
Choose $\delta$ and then $S_0,\rho_0$ so that the total subtracted bracket
is at most $\kappa$.  This proves \eqref{JP:eq:FS}.  Finally
\eqref{JP:eq:exact-flux-forcing} gives the displayed supremum identity.
\end{proof}

\begin{proposition}[Archived exact reduction of JP--BR]
\label{JP:prop:BR-reduction}
JP--BU and JP--FS give the former JP--BR interface, namely
\eqref{JP:eq:fan-upper} and
\begin{equation}\label{JP:eq:BR}
 \Phi_\alpha(h_\alpha\mathbf1)
 -\Phi_{\alpha_*}(h_{\alpha_*}\mathbf1)
 -\frac1{2c}\|\ell_\alpha\|_{(\Gjp_\alpha)^*}^2
 \ge c_0\Jfour(\alpha).
\end{equation}
\end{proposition}

\begin{proof}
Equation \eqref{JP:eq:exact-flux-forcing} and
$\mathcal Q_\alpha=c\Gjp_\alpha$ give the exact identity
\[
 \frac1{2c}\|\ell_\alpha\|_{(\Gjp_\alpha)^*}^2
 =\frac12\sup_{z\ne0}
 \frac{(\sum_i r_{\alpha,i}z_i)^2}{\mathcal Q_\alpha(z)}.
\]
Now apply JP--BU and JP--FS.
\end{proof}

\begin{remark}[Why the closed DM row did not save JE--FM]
In disk coordinates the first nonzero part of the exact covector is
\[
 2q_D\bigl(f_\alpha,
       v_\alpha-w_\alpha f_\alpha'\bigr).
\]
DM controls the null-form contribution in the $v_\alpha$ slot, but not the
curvature placement $-2q_D(f_\alpha,w_\alpha f_\alpha')$.  The certified
three-periodic CT family makes this curvature placement comparable to
$a$, while
$\Jfour(\alpha)^{1/2}\Gjp_\alpha(z)^{1/2}$ is also comparable to $a$.
Thus the archived material model does not behave as an $o(1)$ perturbation
of the disk Schur margin.  The old stationary O--H determinant cannot be
inserted at the nonstationary tangential polygon either.  After the
corner-transfer retraction in
Remark~\ref{JP:rem:period-three-material-forcing}, this is only a diagnostic,
not a proof that the genuine separate covector estimate is false.  The
fixed-loss scalar JP--JC, rather than an unproved citation to TR or DM, is
the active replacement.
\end{remark}

\begin{theorem}[Archived separated completion]
\label{JP:thm:completion}
Assume JP--SC and JP--JE.  Let $\mathcal R_\alpha$ be the Riesz map of
$\Gjp_\alpha$ and set
\begin{equation}\label{JP:eq:corrector}
 z_\alpha^\sharp=-c^{-1}\mathcal R_\alpha^{-1}\ell_\alpha.
\end{equation}
Then, for every normalized
$P_\alpha(h_0+z)\in\mathcal U_{S_0,\rho_0}$,
\begin{align}
 \Gjp_\alpha(z_\alpha^\sharp)&\le C\Jfour(\alpha),
 \label{JP:eq:corrector-bound}\\
 \Phi_\alpha(h_\alpha\mathbf1+z)
 -\Phi_{\alpha_*}(h_{\alpha_*}\mathbf1)
 &\ge c_0\Jfour(\alpha)
 +\frac c2\Gjp_\alpha(z-z_\alpha^\sharp).
 \label{JP:eq:joint-gap}
\end{align}
Thus JP--SC plus JP--JE implies the active-chart coercive node directly;
there is no support-direction CT--M4 or post-cubic remainder.
\end{theorem}

\begin{proof}
Proposition~\ref{JP:prop:FS} supplies JP--FS from the two hypotheses; together
with the closed JP--BU theorem, Proposition~\ref{JP:prop:BR-reduction} supplies
the baseline Schur margin.
By JP--SC and \eqref{JP:eq:exact-integration},
\[
 \Phi_\alpha(h_\alpha\mathbf1+z)-\Phi_\alpha(h_\alpha\mathbf1)
 \ge \ell_\alpha[z]+\frac c2\Gjp_\alpha(z).
\]
Hilbert completion gives
\[
 \ell_\alpha[z]+\frac c2\Gjp_\alpha(z)
 =\frac c2\Gjp_\alpha(z-z_\alpha^\sharp)
 -\frac1{2c}\|\ell_\alpha\|_{(\Gjp_\alpha)^*}^2.
\]
Proposition~\ref{JP:prop:BR-reduction} supplies \eqref{JP:eq:BR}; add it to
obtain \eqref{JP:eq:joint-gap}.  Equations
\eqref{JP:eq:fan-upper} and \eqref{JP:eq:BR} imply
$\|\ell_\alpha\|_{(\Gjp_\alpha)^*}^2\le C\Jfour(\alpha)$; use
\eqref{JP:eq:corrector} to obtain \eqref{JP:eq:corrector-bound}.
\end{proof}

\subsection{Audit disposition}

\begin{center}
\begin{longtable}{>{\raggedright\arraybackslash}p{.34\linewidth}
                  >{\raggedright\arraybackslash}p{.54\linewidth}}
\toprule
item & disposition\\
\midrule
fixed-normal vertex/support path & exact affine identity
\eqref{JP:eq:vertex-linear};\\
genuine side-normal velocity & constant $z_i$, equation
\eqref{JP:eq:constant-normal};\\
JP--EH & closed by Theorem~\ref{JP:thm:EH};\\
JP--NS & closed by Proposition~\ref{JP:prop:normal-speed-chart-entry}; the
downstream metric is fixed by \eqref{JP:eq:JP-metric-fixed};\\
raw $q_D$ norm of the piecewise-constant true normal & rejected as
ill-typed because vertex jumps have logarithmically divergent
$H^{1/2}$ energy;\\
disk cubic candidate & rigorously noncoercive on the alternating family of
Proposition~\ref{JP:prop:cubic-counterexample};\\
JP--SC & archived stronger route; estimates \eqref{JP:eq:S0} and
\eqref{JP:eq:ST} are not active premises;\\
JP--BU & closed by Theorem~\ref{JP:thm:BU};\\
exact support forcing & centered side-flux covector
\eqref{JP:eq:exact-flux-forcing}--\eqref{JP:eq:flux-moments};\\
$D0_{\rm ns}$ and JE--P & closed by
Theorem~\ref{JP:thm:normal-disk-endpoint};\\
JE--V & closed by Theorem~\ref{JP:thm:normal-value-matching};\\
JE--FM & unavailable; its former period-three material falsification is
retracted by Remark~\ref{JP:rem:period-three-material-forcing}, and the row is
absent from the active DAG;\\
vanishing joint comparison & former material counterexample retracted by
Remark~\ref{JP:rem:period-three-material-joint-loss}; the active route uses a
fixed absorbable loss in JP--JC and does not assume this statement;\\
JP--FS and JP--BR & archived conditional algebra; unnecessary on the
JP--JC route;\\
support post-cubic and support-linear fourth tails & removed from the DAG
by the exact value identity \eqref{JP:eq:exact-integration}.\\
\bottomrule
\end{longtable}
\end{center}
\endgroup
